\section{The target and the scope of this milestone}

Reisner's criterion converts an algebraic property of a face ring into
vanishing of the homology of every link. It makes Cohen--Macaulayness
accessible to combinatorial topology and explains why triangulated spheres
give such rings. Our target is precisely Panov, Theorem 2.5
\cite{panov}, over an arbitrary field. The empty face is essential.
This document gives a complete informal proof route. Only the nodes explicitly
marked as checked have Lean implementations; all other nodes, including
foundational comparison theorems, are future formalization work.

The route is through the multigraded extended Cech complex, local cohomology,
and algebraic depth. Panov's Theorem 2.9 is a formula for Tor using induced
subcomplexes. It is not the local-cohomology formula using links below.
We compute the latter directly, including its degree shift and signs.
The proofs of the project-specific lemmas below are arguments supplied by this
project; the externally used algebraic foundations are separately identified.

\section{Objects and conventions}

\begin{definition}[Augmented finite complexes]\label{def:complex}
\lean{Reisner.Complex, Reisner.Complex.face_of_subset, Reisner.Complex.nonface_of_subset}
\leanok
Fix a finite type $V$, possibly empty. A complex $K$ on the ambient set $V$
is a downward closed collection of finite subsets of $V$ containing
$\varnothing$. The complex $\{\varnothing\}$ is allowed; the void collection
is excluded. The Lean structure also supports infinite ambient types.
\end{definition}

\begin{definition}[Vertices, facets, and dimension]\label{def:dimension}
\uses{def:complex}
\lean{Reisner.Complex.vertices, Reisner.Complex.maxFaceCard, Reisner.Complex.card_le_maxFaceCard, Reisner.Complex.exists_face_card_eq_maxFaceCard, Reisner.Complex.exists_facet_card_eq_maxFaceCard, Reisner.Complex.dimension, Reisner.Complex.maxFaceCard_eq_zero_iff_no_vertices, Reisner.Complex.dimension_eq_neg_one_iff_no_vertices}
\leanok
The actual vertices of $K$ are
$W=\{v\in V:\{v\}\in K\}$; elements of $V\setminus W$ are ghost vertices.
A facet is an inclusion-maximal face. Put
\[
 d(K)=\max_{F\in K}|F|,\qquad \dim K=d(K)-1\in\mathbb Z.
\]
Thus $\dim\{\varnothing\}=-1$, not zero.
In Lean, $d(K)$ is \texttt{maxFaceCard}, the finite supremum of the
cardinalities of all faces, and \texttt{dimension} is its integer cast
minus one. The empty face makes the face type nonempty, so a face attains
this maximum; extension to a facet shows that a facet also attains it.
The checked zero-cardinality and dimension-minus-one characterizations
hold exactly when all ambient vertices are ghosts.
\end{definition}

The Lean representation is a structure whose faces form a
\texttt{LowerSet (Finset V)}, with an explicit empty-face field. This reuses
mathlib's order API. The mathlib types
\[
\texttt{PreAbstractSimplicialComplex},\qquad\texttt{AbstractSimplicialComplex}
\]
use different conventions: the first excludes the empty face and the second
also requires every singleton. Removing the empty face
gives the former representation; adjoining it is the inverse operation.
No such conversion is needed for this milestone. Finiteness of $V$ is a
hypothesis of the ultimate target, but is unnecessary for monomial membership
and the standard monomial basis.

\begin{definition}[Links]\label{def:link}
\uses{def:complex,def:dimension}
\lean{Reisner.Complex.link, Reisner.Complex.mem_link_iff, Reisner.Complex.link_empty, Reisner.Complex.notMem_link_singleton_of_mem, Reisner.Complex.mem_link_facet_iff}
\leanok
For a face $\sigma\in K$ define a complex on the same ambient set by
\[
 \operatorname{lk}_K\sigma
 =\{\tau:\tau\cap\sigma=\varnothing,\ \sigma\cup\tau\in K\}.
\]
Its empty face is present. Vertices of $\sigma$ become ghosts.
In particular $\operatorname{lk}_K\varnothing=K$, and the link of a facet
is $\{\varnothing\}$. These statements hold for an arbitrary ambient set
$V$, with all faces finite.
\end{definition}

\begin{lemma}[Link calculus]\label{lem:link-calculus}
\uses{def:link,def:dimension}
If $\tau\in\operatorname{lk}_K\sigma$, then
$\operatorname{lk}_{\operatorname{lk}_K\sigma}\tau
=\operatorname{lk}_K(\sigma\cup\tau)$.
Facets of $\operatorname{lk}_K\sigma$ correspond to facets of $K$ containing
$\sigma$ by $G\mapsto G\setminus\sigma$.
Always $\dim\operatorname{lk}_K\sigma\le d(K)-|\sigma|-1$.
If every facet of $K$ has cardinality $d$, equality holds with $d(K)=d$.
The link of a pure finite complex is pure as well. The same-ambient
formulation retains ghost vertices.
\lean{Reisner.Complex.link_link, Reisner.Complex.mem_facets_link_iff, Reisner.Complex.linkFacetEquiv, Reisner.Complex.link_face_card_add_le_maxFaceCard, Reisner.Complex.link_maxFaceCard_add_card_le, Reisner.Complex.link_dimension_le, Reisner.Complex.link_maxFaceCard_eq_zero_of_facet, Reisner.Complex.link_dimension_eq_neg_one_of_facet, Reisner.Complex.link_maxFaceCard_add_card_eq_of_pure, Reisner.Complex.link_dimension_eq_of_pure, Reisner.Complex.link_pure_of_pure}
\leanok
\end{lemma}
\begin{proof}
Expand the two disjointness conditions in an iterated link: a face $T$ is
disjoint from both $\sigma$ and $\tau$ and its union with both is in $K$.
This is exactly membership in the stated link. For the facet correspondence,
strict enlargement on either side gives a strict enlargement on the other by
union or difference with $\sigma$. A maximal extension exists because $V$ is
finite. The disjoint union $\sigma\cup T$ has size $|\sigma|+|T|\le d(K)$.
For a pure complex, extend $\sigma$ to a facet and apply the correspondence;
all resulting facets have size $d-|\sigma|$.
\end{proof}

\begin{definition}[Reduced homology, including degree minus one]\label{def:homology}
\uses{def:complex}
Choose a total order on $V$. For a field $k$, the augmented chain complex has
$\widetilde C_i(K;k)$ freely spanned by increasing faces of size $i+1$ for
$i\ge0$, $\widetilde C_{-1}=k$ with empty-face basis, and zero terms for
$i<-1$. The boundary deletes vertices with signs $(-1)^r$; in degree zero
each vertex maps to $1\in k$. Define $\widetilde H_i$ as kernel modulo image.
Define reduced cochains by $\widetilde C^i=\operatorname{Hom}_k(\widetilde C_i,k)$
with the transposed differential, and write $\widetilde H^i$ for cohomology.
Changing the order gives the usual orientation-sign chain isomorphism.
In particular $\widetilde H_{-1}(\{\varnothing\};k)=k$, whereas
$\widetilde H_{-1}(K;k)=0$ if $K$ has any actual vertex. All groups below
degree $-1$ vanish.
\end{definition}
The boundary squares to zero: deleting two different vertices in either
order gives the same remaining face and opposite signs, and the boundary
of an edge augments to $1-1=0$. Relabelling the order multiplies an oriented
face by the sign of its sorting permutation, which commutes with deletion.
The checked Lean chain model is the dual of the actual-face cochain model,
with the signed boundary and the empty-face degree. Order-comparison
isomorphisms and the canonical homology-cohomology isomorphism remain work.

\begin{lemma}[The augmented face-cardinality cochain complex]\label{lem:augmented-cochain-complex}
\uses{def:complex}
\lean{Reisner.Complex.SimplicialFaceIndex, Reisner.Complex.SimplicialCochain, Reisner.Complex.SimplicialCochainDirectSum, Reisner.Complex.simplicialCochainDirectSumEquiv, Reisner.Complex.simplicialCochainBasis, Reisner.Complex.simplicialCochainBasis_apply, Reisner.Complex.simplicialFaceEraseIndex, Reisner.Complex.simplicialRankBelow, Reisner.Complex.simplicialDifferential, Reisner.Complex.simplicialDifferential_apply, Reisner.Complex.simplicialDifferential_comp_eq_zero, Reisner.Complex.simplicialFaceIndexZeroUnique, Reisner.Complex.simplicialCochainZeroEquiv, Reisner.Complex.simplicialCochainZeroEquiv_apply, Reisner.Complex.simplicialDifferential_zero_apply, Reisner.Complex.simplicialDifferential_zero_of_no_vertices, Reisner.Complex.simplicialDifferential_basis_apply, Reisner.Complex.simplicialCochain_subsingleton_of_card_lt, Reisner.Complex.simplicialCochain_isZero_of_card_lt, Reisner.Complex.simplicialCochainComplex, Reisner.Complex.simplicialCochainComplex_X, Reisner.Complex.simplicialCochainComplex_d}
\leanok
Let $V$ be a finite linearly ordered type, $K$ an augmented complex on
$V$, and $R$ any commutative ring. For $n\in\mathbb N$, set
\[
 \mathcal F_n(K)=\{F\in K:|F|=n\},\qquad
 E^n(K;R)=\prod_{F\in\mathcal F_n(K)}R.
\]
The product is finite and has the coordinate basis $e_F$, whose value is
$1$ at $F$ and $0$ at all other faces. It is canonically the corresponding
finite direct sum. If $G\in\mathcal F_{n+1}(K)$ and $v\in G$, its deletion
$G\setminus\{v\}$ belongs to $\mathcal F_n(K)$. Write
$r_G(v)=|\{u\in G:u<v\}|$. The linear maps
\[
 (\delta^n f)(G)=\sum_{v\in G}(-1)^{r_G(v)}
                         f(G\setminus\{v\})
\]
satisfy $\delta^{n+1}\delta^n=0$ for every $n\geq0$ and give an actual
cochain complex $E^\bullet(K;R)$ in the category of $R$-modules. Its term
at $n$ and its differential from $n$ to $n+1$ are exactly the module and
map displayed above. On the coordinate basis the following coordinate
formula holds:
\[
 (\delta^n e_F)(G)=\sum_{v\in G}(-1)^{r_G(v)}
       \begin{cases}
        1,&G\setminus\{v\}=F,\\
        0,&G\setminus\{v\}\ne F.
       \end{cases}
\]
The face-cardinality index $n$ represents reduced cochain degree
$(n:\mathbb Z)-1$; in particular $n=0$ represents degree $-1$.
There is no incoming differential at cardinality index zero. Evaluation at the
unique empty face gives a canonical linear equivalence $E^0(K;R)\simeq_R R$.
Under this equivalence, $\delta^0$ sends a scalar $r$ to the constant
cochain with value $r$ on the actual singleton faces. Ghost singletons
are absent from this index set. If every singleton is a nonface, the linear
map $\delta^0$ is zero. For $n>|V|$, the actual module $E^n(K;R)$ is a
zero object in the category of $R$-modules. These assertions include empty
$V$, ghost vertices, characteristic two, and the zero coefficient ring.
\end{lemma}
\begin{proof}
\leanok
Faces of a fixed cardinality form a finite subtype of the finite subsets
of $V$, so the coordinate product and its standard basis are defined.
Downward closure puts every deletion back in $K$, and deletion of a vertex
of an $(n+1)$-element set leaves exactly $n$ vertices. Coordinate evaluation,
scalar multiplication by a sign, and finite summation are linear, proving
linearity of $\delta^n$.

Expand $\delta^{n+1}\delta^n f$ at a face $H$ of size $n+2$.
Its terms are indexed by ordered pairs of distinct vertices $v,w\in H$,
with coefficient
\[
 (-1)^{r_H(v)+r_{H\setminus\{v\}}(w)}
                      f(H\setminus\{v,w\}).
\]
Swapping $v$ and $w$ is a fixed-point-free involution on these indices.
The two deletions give the same remaining face. If $v<w$, deleting $v$
lowers the rank of $w$ by one, whereas deleting $w$ leaves the rank of
$v$ unchanged; the case $w<v$ is symmetric. The two sign products are
therefore additive inverses, and every pair cancels. This proves the
square-zero identity also for $n=0$, where the two terms at an edge are
the values at its two singleton faces with opposite signs. Cancellation
uses additive inverses, with no division by two and no restriction on the
characteristic. The square-zero maps construct the genuine module-valued
cochain complex.

Substituting the coordinate value of $e_F$ into the deletion formula gives
the stated basis-action formula. The cardinality-zero face is uniquely
$\varnothing$ by the augmented convention, so evaluation there gives the
scalar equivalence. Every singleton has deletion $\varnothing$ and deletion
rank zero, giving the constant value in the dual augmentation. If there
are no actual vertices, there are no singleton indices, and the outgoing
linear map is zero. No ghost contributes an index. Finally, if $n>|V|$,
there is no $n$-element subset of $V$, so the index type is empty and its
function module is a zero object. All constructions and cancellations work
unchanged for the zero ring.
\end{proof}

\begin{lemma}[Coordinate formula for basis insertion]\label{lem:augmented-basis-insertion}
\uses{lem:augmented-cochain-complex}
For an $n$-face $F$ and an $(n+1)$-face $G$, the coordinate of the
differential of $e_F$ is
\[
 (\delta^n e_F)(G)=
 \sum_{v\in G}
   \begin{cases}
    (-1)^{r_F(v)},&v\notin F\text{ and }F\cup\{v\}=G,\\
    0,&\text{otherwise}.
   \end{cases}
\]
Thus the coefficient is nonzero only when $G$ is obtained by inserting one
vertex into $F$, and the sign is the number of vertices of $F$ below it.
\lean{Reisner.Complex.simplicialDifferential_basis_insertion_apply}
\leanok
\end{lemma}
\begin{proof}
In the deletion-coordinate formula, the basis cochain $e_F$ contributes at
the deletion of $v$ precisely when $G\setminus\{v\}=F$. This is equivalent
to $v\notin F$ and $F\cup\{v\}=G$. Under this condition the vertices below
$v$ are unchanged by deleting $v$, so $r_G(v)=r_F(v)$. The Lean theorem proves
this coordinate formula directly.
\end{proof}

\begin{lemma}[Degree-minus-one cohomology endpoints]\label{lem:augmented-cohomology-minus-one}
\uses{lem:augmented-cochain-complex}
The actual homology object at cardinality index zero of the constructed
cochain complex is canonically isomorphic to $\ker\delta^0$. If $K$ has no
actual vertices, equivalently $K=\{\varnothing\}$, it is canonically
linearly equivalent to $R$. If $K$ has an actual vertex, it is zero.
These assertions include empty $V$, ghost vertices, characteristic two,
and the zero coefficient ring.
\lean{Reisner.Complex.simplicialCochainHomologyZeroIsoKer, Reisner.Complex.simplicialCochainHomologyZeroEquivOfNoVertices, Reisner.Complex.simplicialCochainHomologyZero_isZero_of_vertex}
\leanok
\end{lemma}
\begin{proof}
Mathlib identifies degree-zero homology of a cochain complex indexed by
$\mathbb N$ with its cycles object, since there is no incoming boundary.
The cycles object is the categorical kernel of $\delta^0$, and in
$\mathrm{ModuleCat}(R)$ this is the usual linear-map kernel. If there are
no actual vertices, $\delta^0=0$ and its kernel is the full empty-face term
$R$. If $\{v\}$ is an actual face, evaluation at that singleton recovers
the empty-face scalar, so $\delta^0$ is injective and the kernel is zero.
The Lean declarations construct the actual homology-to-kernel isomorphism
and prove both endpoint cases.
\end{proof}

\begin{lemma}[Integer-indexed augmented cochains and grading endpoints]\label{lem:augmented-cochains}
\uses{lem:augmented-cochain-complex,lem:augmented-cohomology-minus-one}
\lean{Reisner.Complex.simplicialCochainDegreeEmbedding, Reisner.Complex.simplicialCochainComplexInt, Reisner.Complex.simplicialCochainComplexIntXIso, Reisner.Complex.simplicialCochainComplexIntRestrictionIso, Reisner.Complex.simplicialCochainComplexIntHomologyIso, Reisner.Complex.simplicialCochainComplexInt_d, Reisner.Complex.simplicialCochainComplexInt_isZero_of_lt_minus_one, Reisner.Complex.simplicialCochainComplexInt_d_eq_zero_of_lt_minus_one, Reisner.Complex.simplicialCochainComplexInt_homology_isZero_of_lt_minus_one, Reisner.Complex.simplicialCochainComplex_homology_isZero_of_card_lt, Reisner.Complex.simplicialCochainInt_natDegree, Reisner.Complex.simplicialCochainInt_natBound, Reisner.Complex.simplicialCochainComplexInt_isZero_of_card_lt, Reisner.Complex.simplicialCochainComplexInt_homology_isZero_of_card_lt, Reisner.Complex.simplicialCochainComplexIntHomologyMinusOneEquivOfNoVertices, Reisner.Complex.simplicialCochainComplexIntHomologyMinusOne_isZero_of_vertex}
\leanok
Let $V$, $K$, and $R$ be as in
Lemma~\ref{lem:augmented-cochain-complex}, and write $E^\bullet$ for its
checked cochain complex indexed by natural face cardinalities. Set
\[
 e(n)=(n:\mathbb Z)-1,\qquad n\in\mathbb N.
\]
There is a genuine cochain complex $C^\bullet(K;R)$ indexed by
$\mathbb Z$, obtained by extending $E^\bullet$ by zero along $e$.
For $i\geq-1$, put $n_i=(i+1).\mathrm{toNat}$, where conversion to a
natural number is used only after verifying $i+1\geq0$. Thus
$(n_i:\mathbb Z)=i+1$ and $e(n_i)=i$. Its terms have canonical
identifications
\[
 C^i(K;R)\cong
 \begin{cases}
  E^{n_i}(K;R),&i\geq-1,\\
  0,&i<-1.
 \end{cases}
\]
Write $\phi_n:C^{e(n)}\simeq_R E^n$ for the term identifications.
Since $e(n+1)=e(n)+1$, the differential satisfies
\[
 \phi_{n+1}\circ d_C^{e(n)}\circ\phi_n^{-1}=\delta^n.
\]
Every differential with source below degree $-1$ is zero. In particular
the map from degree $-1$ to degree $0$ is the checked dual augmentation,
and the incoming map from degree $-2$ to degree $-1$ is zero.
Restriction of $C^\bullet$ along $e$ is isomorphic to $E^\bullet$ as an
actual cochain complex; the isomorphism and its inverse have components
$\phi_n$ and $\phi_n^{-1}$. There are canonical isomorphisms of actual
cohomology objects
\[
 H^{e(n)}(C^\bullet)\cong H^n(E^\bullet),\qquad n\in\mathbb N.
\]

For every $i<-1$, both the module $C^i$ and its actual cohomology object
$H^i(C^\bullet)$ are zero. For every
$i>(|V|:\mathbb Z)-1$, they are again both zero; this upper inequality
is equivalent to $|V|<n_i$ when $i\geq-1$. In particular the actual
homology object $H^n(E^\bullet)$ is zero for every natural $n>|V|$.
At degree $-1$ the displayed cohomology isomorphism identifies
$H^{-1}(C^\bullet)$ with the checked cardinality-zero cohomology:
it is canonically linearly equivalent to $R$ if $K$ has no actual
vertices, equivalently $K=\{\varnothing\}$, and is zero if an actual
vertex exists. All statements include empty $V$, ghosts, characteristic
two, and the zero coefficient ring. No comparison with the Cech complex
or local cohomology is asserted here.
\end{lemma}
\begin{proof}
\leanok
The map $e$ is injective, preserves successive degrees, and has image
exactly the integers at least $-1$. Extend the checked complex along this
embedding of cochain shapes. On its image the extension has the same
objects and maps up to the canonical transport isomorphisms; outside its
image it has zero objects and zero maps. Its square-zero law is inherited
from $E^\bullet$ when both successive maps are in the image, and otherwise
one of the maps is zero. This includes the incoming map at degree $-1$.
The term transports commute with every differential, so they and their
inverses define the claimed restriction isomorphism.

For $i\geq-1$, nonnegativity of $i+1$ gives
$((i+1).\mathrm{toNat}:\mathbb Z)=i+1$, hence $e(n_i)=i$.
Moreover $i>(|V|:\mathbb Z)-1$ is equivalent to $|V|<n_i$:
add one in $\mathbb Z$ and use the same cast identity. These identities
justify the natural indices without truncated subtraction. The extension
identifies the kernels and images determining cohomology at every image
degree with those of $E^\bullet$. At $e(0)=-1$, the incoming term of the
extension is zero, matching the absence of an incoming differential at
natural index zero, so this comparison includes that endpoint.

Below degree $-1$, the term itself is zero. Above the stated upper bound,
the corresponding natural index exceeds $|V|$, and
Lemma~\ref{lem:augmented-cochain-complex} makes that term zero as well.
A complex is exact at a zero term, so its actual cohomology object is
zero there. The same argument proves high-cardinality vanishing directly
for $E^\bullet$. Finally transport the scalar equivalence and the
actual-vertex vanishing from
Lemma~\ref{lem:augmented-cohomology-minus-one} through the cohomology
isomorphism at $n=0$. No coefficient is divided out, so the argument
includes the zero ring and characteristic two.

The extension and its term, restriction, and homology isomorphisms are
implemented using mathlib's embedding API. The chain model,
order-independence comparison, and the full reduced homology and duality assertions of
Definition~\ref{def:homology} remain separate formalization work.
\end{proof}

\begin{lemma}[Finite-dimensional augmented cochain terms]\label{lem:cochain-finite-dim}
\uses{lem:augmented-cochains}
Let $V$ be finite and $k$ a field. The augmented cochain term in
face-cardinality degree $n$ has dimension equal to the number of actual
$n$-element faces of $K$. In degree $0$, this counts the empty face; ghost
vertices contribute no basis coordinates.
\lean{Reisner.Complex.simplicialCochain_finiteDimensional, Reisner.Complex.simplicialCochain_finrank}
\leanok
\end{lemma}
\begin{proof}
The coordinate basis is indexed by actual faces of cardinality $n$. Its index
type is finite because the ambient vertex type is finite. The finite-basis
dimension theorem gives both finite dimensionality and the stated rank.
\end{proof}

\begin{lemma}[Augmented simplicial chain complex]\label{lem:augmented-chain-complex}
\uses{lem:cochain-finite-dim,lem:augmented-cochain-complex}
For a finite ambient vertex type and a field $k$, take the degree-$n$ chain
space to be the dual of the cochain space indexed by actual $n$-element faces.
It is finite-dimensional with one basis vector per actual face. Evaluation
gives a canonical isomorphism from cochains to the dual of chains, compatible
with the cochain differential and dual chain boundary. The dual differential
gives the boundary
\[
 \partial_n : C_n(K;k)\longrightarrow C_{n-1}(K;k),
 \qquad C_n(K;k)=\widetilde C_{n-1}(K;k),
\]
and $\partial_{n-1}\partial_n=0$. On a basis face it is the signed sum of
its codimension-one faces. The empty face is the basis in face-cardinality
degree zero (reduced degree $-1$); ghost vertices are absent. Reindexing
puts face-cardinality degree $n$ in integer reduced degree $n-1$ and extends
by zero below $-1$.
\lean{Reisner.Complex.simplicialChainComplex, Reisner.Complex.simplicialChain_finiteDimensional, Reisner.Complex.simplicialChain_finrank, Reisner.Complex.simplicialCochainEvalEquivDualChain, Reisner.Complex.simplicialCochainEvalEquivDualChain_naturality, Reisner.Complex.simplicialBoundary_comp_eq_zero, Reisner.Complex.simplicialBoundary_basis_apply, Reisner.Complex.simplicialChainComplexInt, Reisner.Complex.simplicialChainComplexIntHomologyIso, Reisner.Complex.simplicialChainComplexInt_homology_isZero_of_lt_minus_one}
\leanok
\end{lemma}
\begin{proof}
The cochain terms have finite coordinate bases indexed by actual faces. Their
dual bases give finite-dimensional chain spaces. Bidual evaluation gives the
stated canonical term isomorphism, and naturality of evaluation proves its
differential compatibility. The chain differential is the linear dual of
cochain deletion, so its basis evaluation is the signed deletion formula.
Dualizing the cochain square-zero identity proves the chain square-zero law.
\end{proof}

\begin{lemma}[Vanishing part of homology-cohomology duality]\label{lem:dual-vanishing}
\uses{lem:augmented-chain-complex}
For each face-cardinality degree $n\in\mathbb N$, the augmented cochain
homology object is zero if and only if the corresponding augmented chain
homology object is zero. This is only the vanishing equivalence; it does not
assert the canonical isomorphism of homology groups in
Lemma~\ref{lem:dual}.
\lean{Reisner.Complex.simplicialCochain_homology_isZero_iff_chain_homology_isZero}
\leanok
\end{lemma}
\begin{proof}
At positive face-cardinality degrees, exactness says that the incoming
differential has range equal to the outgoing differential's kernel. Taking
linear duals exchanges ranges with annihilators of kernels and kernels with
annihilators of ranges. At degree zero, the incoming cochain map is zero,
while chain exactness says the dual of the augmentation differential is
surjective; over a field this is equivalent to injectivity of augmentation.
\end{proof}

\begin{lemma}[Canonical homology--cohomology duality in face-cardinality degrees]\label{lem:dual-isomorphism}
\uses{lem:augmented-chain-complex}
For every $n\in\mathbb N$, finite ordered vertex set $V$, augmented
complex $K$ on $V$, and field $k$, evaluation induces a canonical linear
isomorphism
\[
 H_n(E^\bullet(K;k))\cong
 \operatorname{Hom}_k(H_n(C_\bullet(K;k)),k),
\]
where the face-cardinality index $n$ represents reduced degree $n-1$.
The statement includes $n=0$, the empty-face degree. Its integer-indexed
extension is Lemma~\ref{lem:dual}.
\lean{Reisner.DualHomology.homologyDualEquiv, Reisner.DualHomology.homologyDualEquiv_apply_cycle, Reisner.Complex.simplicialCohomologyDualHomologyEquiv}
\leanok
\end{lemma}
\begin{proof}
For a short complex of vector spaces, restriction of a dual cycle to the
original cycles annihilates boundaries, so it descends to homology. Every
functional on homology extends to the cycle space and then to the middle
space, giving surjectivity; the kernel consists exactly of dual boundaries.
This yields the canonical evaluation isomorphism. The finite face-coordinate
duality commutes with the signed differentials, and the short complexes
around degree $n$ transport the result to the actual simplicial complexes.
\end{proof}

\begin{lemma}[Homology versus cohomology]\label{lem:dual}
\uses{lem:dual-isomorphism,lem:augmented-cochains,lem:augmented-chain-complex}
Let $K$ be a finite augmented complex on an ordered ambient vertex set and
$k$ a field. For every integer $i$, canonical evaluation gives
\[
\widetilde H^i(K;k)\cong\operatorname{Hom}_k(\widetilde H_i(K;k),k).
\]
Consequently their vanishing conditions agree. Both terms and both homology
objects vanish for $i<-1$ and for $i>\dim K=d(K)-1$; in particular they
vanish for $i>|V|-1$. At $i=-1$, the existing augmented empty-face convention
is preserved: cohomology is $k$ when $K=\{\varnothing\}$, even with ghost
vertices, and is zero when $K$ has an actual vertex. The same vanishing
alternative holds for homology.
\lean{Reisner.Complex.simplicialCohomologyIntDualHomologyEquiv, Reisner.Complex.simplicialCohomologyIntDualHomologyEquiv_reindex, Reisner.Complex.simplicialCochainInt_homology_isZero_iff_chainInt_homology_isZero, Reisner.Complex.simplicialCochainComplexInt_isZero_of_dimension_lt, Reisner.Complex.simplicialCochainComplexInt_homology_isZero_of_dimension_lt, Reisner.Complex.simplicialChainComplexInt_homology_isZero_of_lt_minus_one, Reisner.Complex.simplicialChainComplexInt_isZero_of_card_lt, Reisner.Complex.simplicialChainComplexInt_homology_isZero_of_card_lt, Reisner.Complex.simplicialChainComplexInt_isZero_of_dimension_lt, Reisner.Complex.simplicialChainComplexInt_homology_isZero_of_dimension_lt}
\leanok
\end{lemma}
\begin{proof}
For $i\geq-1$, put $n=(i+1).\mathrm{toNat}$. Then $(n:\mathbb Z)-1=i$.
Use the canonical extension homology isomorphisms for chains and cochains,
apply Lemma~\ref{lem:dual-isomorphism} at $n$, and precompose functionals
with the chain homology isomorphism. This includes $n=0$, so the augmented
degree $-1$ is transported without discarding the empty face. For $i<-1$
both homology objects are zero and the unique linear equivalence is canonical.
A vector space has zero linear dual exactly when it is zero, giving the
vanishing equivalence in all integer degrees.

For integer $i>\dim K=d(K)-1$, the corresponding face cardinality
$i+1$ exceeds $d(K)$, so there are no actual faces in that term.
Both the cochain term and its dual chain term are zero. Transport along
the same degree embedding proves that both integer complexes are exact
there. The ambient-cardinality bounds follow in the same way. The degree
$-1$ alternatives are those already proved for the cochain complex, together
with the vanishing equivalence. Ghost vertices add no actual face coordinates.
\end{proof}

\begin{lemma}[Reduced degree zero and graph connectivity]\label{lem:homology-connectivity}
\uses{def:homology,lem:augmented-chain-complex,lem:dual}
For a finite augmented complex over a field $k$ with at least one actual vertex,
$\widetilde H_0(K;k)=0$ if and only if its graph of actual vertices and
edges is connected. This graph comparison is not yet formalized.
\end{lemma}
\begin{proof}
Boundaries identify the basis vectors at the endpoints of each edge.
Quotienting by these relations leaves one basis vector per graph component.
The augmentation is the sum of their coordinates, whose kernel has dimension
one less than the number of components. This is a separate comparison with
the graph model, not a hypothesis of the checked duality theorem.
\end{proof}

\begin{definition}[Nonface generators and the face ring]\label{def:ring}
\uses{def:complex}
\lean{Reisner.squarefreeExponent, Reisner.monomial_squarefreeExponent_eq_prod, Reisner.Complex.nonfaceExponents, Reisner.Complex.stanleyReisnerIdeal, Reisner.Complex.stanleyReisnerIdeal_eq_span, Reisner.Complex.FaceRing}
\leanok
Let $S=k[x_v:v\in V]$, represented by \texttt{MvPolynomial V k}.
Exponents are finitely
supported maps $\alpha:V\to_0\mathbb N$, with
$\operatorname{supp}\alpha=\{v:\alpha(v)\ne0\}$ and monic monomial $x^\alpha$.
For a finite set $F$, let $\mathbf1_F$ be its finitely supported indicator
and $x_F=x^{\mathbf1_F}=\prod_{v\in F}x_v$ (the empty product is $1$).
Define
\[
 I_K=(x_F:F\notin K),\qquad A=k[K]=S/I_K.
\]
The definition is an ideal span of squarefree nonface generators, not a
membership predicate.
\end{definition}
For the future grading development use total degree $1$ for variables, and retain
the finer $\mathbb N^V$ grading. Panov uses degree $2$ instead; this regrading
does not change the ungraded ring or Cohen--Macaulayness.

\begin{definition}[Algebraic Cohen--Macaulayness]\label{def:cm}
For a Noetherian local ring $(B,\mathfrak n)$, a regular sequence is a finite
sequence in $\mathfrak n$ whose successive multiplications on the successive
quotients are injective; the final quotient is nonzero. Its depth is the
supremum of lengths of these sequences. It is Cohen--Macaulay if
$\operatorname{depth} B=\dim B$, with Krull dimension defined by chains of
prime ideals. A Noetherian ring $A$ is Cohen--Macaulay if every localization
$A_{\mathfrak p}$ at a prime is a Cohen--Macaulay local ring.
\lean{Reisner.IsLocalCohenMacaulay, Reisner.IsCohenMacaulay}
\leanok
\end{definition}
These are the algebraic definitions in Stacks, Sections 10.103--10.104
\cite{stacks-cm}. Stacks also gives the equivalent test at all maximal ideals.

\section{The implemented algebraic entry point}

\begin{lemma}[Squarefree exponent comparison]\label{lem:indicator}
\uses{def:ring}
\lean{Reisner.squarefreeExponent_apply, Reisner.support_squarefreeExponent, Reisner.squarefreeExponent_le_iff}
\leanok
For any finite $F$ and $\alpha:V\to_0\mathbb N$,
\[
 \mathbf1_F\le\alpha\quad\Longleftrightarrow\quad
 F\subseteq\operatorname{supp}\alpha.
\]
\end{lemma}
\begin{proof}
\leanok
The order is pointwise. At $v\in F$ it requires $1\le\alpha(v)$, equivalent
over $\mathbb N$ to $\alpha(v)\ne0$. At $v\notin F$ it requires only
$0\le\alpha(v)$. This also shows $\operatorname{supp}\mathbf1_F=F$.
\end{proof}

\begin{theorem}[Monomial membership]\label{thm:membership}
\uses{def:ring,lem:indicator}
\lean{Reisner.Complex.monomial_mem_stanleyReisnerIdeal_iff}
\leanok
For every $\alpha:V\to_0\mathbb N$,
\[
 x^\alpha\in I_K\quad\Longleftrightarrow\quad
 \operatorname{supp}\alpha\notin K.
\]
\end{theorem}
\begin{proof}
\leanok
Mathlib's \texttt{MvPolynomial.mem\_ideal\_span\_monomial\_image} says that a
polynomial is in the ideal spanned by monic monomials with exponent set $E$
exactly when each exponent in its polynomial support dominates an element
of $E$. Apply it with $E=\{\mathbf1_F:F\notin K\}$ to $x^\alpha$.
Its polynomial support is $\{\alpha\}$ because $1\ne0$ in a field.
Thus membership is equivalent to the existence of $F\notin K$ with
$\mathbf1_F\le\alpha$. By Lemma~\ref{lem:indicator} this says that a nonface
is contained in $\operatorname{supp}\alpha$. Downward closure forbids this
when the latter is a face. Conversely, when it is a nonface, take
$F=\operatorname{supp}\alpha$ itself. The argument needs only a nontrivial
commutative semiring, and does not need $V$ finite; the Lean theorem retains
this greater generality.
\end{proof}

\begin{proposition}[Standard monomial basis]\label{prop:basis}
\uses{thm:membership}
\lean{Reisner.Complex.mem_stanleyReisnerIdeal_iff, Reisner.Complex.faceExponents, Reisner.Complex.stanleyReisnerIdeal_restrictScalars_eq_restrictSupport, Reisner.Complex.isCompl_stanleyReisnerIdeal_faceSupport, Reisner.Complex.faceRingBasis, Reisner.Complex.faceRingBasis_apply, Reisner.Complex.faceRing_nontrivial, Reisner.Complex.faceRing_algebraMap_injective}
\leanok
A polynomial $f\in S$ belongs to $I_K$ exactly when every exponent in its
polynomial support has nonface vertex support. The classes of the monomials
$x^\alpha$ with $\operatorname{supp}\alpha\in K$ form a $k$-basis of $A$.
The basis index is the subtype of finitely supported exponents with
face support, and its vector at $\alpha$ is the quotient class of
$x^\alpha$. In particular $A\ne0$ and the scalar map $k\longrightarrow A$
is injective.
\end{proposition}
\begin{proof}
\leanok
Apply the same mathlib monomial-ideal theorem to each exponent in $f$ and
use the indicator comparison and downward closure, exactly as above.
Split the polynomial monomial basis into face-supported and nonface-supported
parts. The ideal is precisely the span of the latter. Every polynomial is a
finite sum of these two parts, so face-supported classes span the quotient.
A linear combination of them in the ideal would have all its nonzero
coefficients simultaneously face-supported and nonface-supported, an
impossibility. Hence they are independent. In Lean, identify the coefficient
submodule of the ideal with \texttt{MvPolynomial.restrictSupport} on the
nonface exponents. The monomial basis and the complementary exponent sets
give complementary submodules. Transport the canonical face-supported basis
using mathlib's \texttt{quotientEquivOfIsCompl} in the inverse direction.
This inverse takes each polynomial
to its class in the actual ideal quotient, as checked explicitly for each
basis vector.
The zero exponent has empty support, a face, so its class $1$ survives.
More generally, a nonzero constant has only the zero exponent in its
polynomial support, so it cannot lie in the ideal. This proves injectivity
of the scalar map. Polynomial membership needs only a commutative semiring;
the basis and scalar injectivity hold over any commutative ring, with
nontriviality of that ring needed for $A\ne0$. The ambient type may be
infinite throughout.
\end{proof}

\begin{lemma}[Weighted homogeneity of the nonface ideal]\label{lem:homogeneous}
\uses{def:ring}
\lean{Reisner.Complex.stanleyReisnerIdeal_isHomogeneous, Reisner.Complex.weightedHomogeneousComponent_mem_stanleyReisnerIdeal, Reisner.weight_fine_eq, Reisner.weight_one_eq_sum}
\leanok
Let $V$ be any type, $K$ an augmented complex on $V$, $R$ a commutative
semiring, $M$ an additive commutative monoid, and $w:V\to M$ any weight
function. For an exponent $\alpha:V\to_0\mathbb N$ put
\[
 |\alpha|_w=\sum_{v\in\operatorname{supp}\alpha}\alpha(v)\,w(v),
 \qquad
 S_d=\left\{p\in R[x_v:v\in V]:
   \forall\alpha,\ [x^\alpha]p\ne0\Longrightarrow |\alpha|_w=d\right\}.
\]
Here $n\,m$ means repeated addition in $M$. These are the weighted
homogeneous submodules of the polynomial algebra; its degree-$d$
projection is
\[
 \pi_d(p)=\sum_{\substack{\alpha\in\operatorname{supp}p\\|\alpha|_w=d}}
       [x^\alpha]p\,x^\alpha.
\]
The existing ideal $I_K=(x_F:F\notin K)$ is homogeneous for this grading:
\[
 p\in I_K\quad\Longrightarrow\quad\pi_d(p)\in I_K
 \qquad(d\in M).
\]
In Lean, \texttt{Ideal.IsHomogeneous} expresses this property for
mathlib's weighted polynomial grading and its component projections.
For $M=V\to_0\mathbb N$ and the fine weight $w(v)=e_v$, with $e_v$ the exponent supported at $v$
and equal to $1$ there, $|\alpha|_w=\alpha$. For the constant weight
$w(v)=1$ in $\mathbb N$, $|\alpha|_w=\sum_v\alpha(v)$, so the same ideal
is homogeneous for both the fine grading and the usual total-degree
grading of the polynomial ring. No finiteness, nonzero-weight, or
nontrivial-coefficient hypothesis is required.
\end{lemma}
\begin{proof}
\leanok
We give the generator argument explicitly; its specialization to the
nonface generating set is a project argument. Every polynomial has only
finitely many exponents, so grouping its terms by $|\alpha|_w$ gives a
finite homogeneous decomposition, even if $V$ is infinite or some weights
are zero. Since $|\alpha+\beta|_w=|\alpha|_w+|\beta|_w$, products of
homogeneous polynomials of degrees $e$ and $f$ are homogeneous of degree
$e+f$. These are the decomposition and multiplication properties of the
weighted polynomial grading.

For each nonface $F$, its generator
$x_F=x^{\mathbf1_F}$ lies in $S_{|\mathbf1_F|_w}$. More generally, a
monomial with any coefficient is homogeneous of its exponent's weighted
degree, including when it is zero. If $q=\sum_e q_e$ is the finite
homogeneous decomposition of an arbitrary multiplier and
$\delta=|\mathbf1_F|_w$, then
\[
 \pi_d(qx_F)=\sum_{\substack{e\\e+\delta=d}}q_e x_F.
\]
Every term on the right is in $I_K$. A polynomial of $I_K$ is a finite
sum of multiples of its displayed generators, and the projections are
additive, so all its projections belong to $I_K$. This proves homogeneity
of the whole generator-defined ideal. Different degrees $e$ may satisfy
$e+\delta=d$; all of them are retained, so no cancellation or ordering
property of $M$ is used. Conversely, if every $\pi_d(p)$ belongs to the
ideal, their finite sum $p$ does also.

In namespace \texttt{Ideal}, the theorem \texttt{homogeneous\_span}
supplies this argument. In namespace \texttt{MvPolynomial}, monomial
homogeneity uses \texttt{isWeightedHomogeneous\_monomial}, and the
polynomial grading is \texttt{weightedGradedAlgebra}. Its decomposition
projections are the stated weighted components, so the library's
component-preservation theorem gives the displayed consequence directly.
Only the set of generators is used: downward closure, the empty-face
witness, and the presence or absence of ambient singleton faces do not
enter this argument. The existing complex conventions are preserved.

For fine weights the degree is
$\sum_v\alpha(v)e_v=\sum_v\alpha(v)\,\mathbf1_{\{v\}}=\alpha$,
the standard reconstruction of a finitely supported function. For unit
weights it is $\sum_v\alpha(v)\cdot1=\sum_v\alpha(v)$, the ordinary
exponent sum. These finite sums also prove the two specializations without
assuming $V$ finite. This lemma concerns the polynomial ideal; construction
of the quotient grading is the separate next proposition.
\end{proof}

\begin{proposition}[Fine grading and degree zero]\label{prop:grading}
\uses{lem:homogeneous,prop:basis,def:ring}
\lean{Reisner.Complex.weightedFaceRingComponent, Reisner.Complex.weightedFaceRingComponent_eq_span, Reisner.Complex.monomial_nonface_eq_zero, Reisner.Complex.weightedFaceRingComponent_eq_polynomial_image, Reisner.Complex.weightedFaceRingComponent_mul, Reisner.Complex.weightedFaceRing_isInternal, Reisner.Complex.weightedFaceRingGradedAlgebra, Reisner.Complex.faceRingFineComponent, Reisner.Complex.faceRingTotalComponent, Reisner.Complex.faceRingFineGradedAlgebra, Reisner.Complex.faceRingTotalGradedAlgebra, Reisner.Complex.faceRingFineComponent_eq_span_monomial, Reisner.Complex.faceRingFineComponent_eq_bot_of_nonface, Reisner.Complex.faceRingTotalComponent_eq_span, Reisner.Complex.faceRingTotalComponent_eq_iSup_fine, Reisner.Complex.faceRingFineComponent_zero, Reisner.Complex.faceRingTotalComponent_zero, Reisner.Complex.faceRingDegreeZeroSubalgebra, Reisner.Complex.faceRingDegreeZeroEquiv, Reisner.Complex.faceRingDegreeZeroEquiv_apply}
\leanok
Allow any ambient type $V$ and any commutative coefficient ring $R$,
in particular any field $k$, and write $A=R[K]$ for the actual quotient.
Let $q:R[x_v:v\in V]\to A$ be its quotient map, let
$J=\{\alpha:V\to_0\mathbb N:\operatorname{supp}\alpha\in K\}$, and let
$b_\alpha=q(x^\alpha)$ for $\alpha\in J$ be the checked standard basis.
For every exponent $\alpha$, the fine component is
\[
 A_\alpha=
 \begin{cases}
 R\,b_\alpha,&\operatorname{supp}\alpha\in K,\\
 0,&\operatorname{supp}\alpha\notin K.
 \end{cases}
\]
These submodules give a genuine internal direct-sum grading of the quotient:
\[
 A=\bigoplus_{\alpha:V\to_0\mathbb N}A_\alpha,
 \qquad 1\in A_0,\qquad A_\alpha A_\beta\subseteq A_{\alpha+\beta}.
\]
Thus the fine grading is indexed by finitely supported nonnegative exponent
vectors; for finite $V$ this is the usual $\mathbb N^V$ grading.
Put $|\alpha|=\sum_v\alpha(v)$, the additive total-degree map. Its
coarsening has components
\[
 T_n=\operatorname{span}_R\{b_\alpha:\alpha\in J,\ |\alpha|=n\}
     =\bigoplus_{|\alpha|=n} A_\alpha.
\]
\[
 A=\bigoplus_{n\in\mathbb N}T_n,
 \qquad T_mT_n\subseteq T_{m+n}.
\]
Both decompositions are supplied as graded-algebra structures on $A$,
with their canonical finite homogeneous decompositions. The fine and total
degree-zero components coincide and are precisely the scalar image:
\[
 A_0=T_0=R\,1=\operatorname{im}(R\longrightarrow A).
\]
With multiplication inherited from $A$, this component is a subalgebra,
and the scalar map induces the canonical algebra equivalence
$R\simeq_R T_0$. In particular the original identification $A_0=k$ holds
over every field. No finite-vertex, purity, no-ghost, nonemptiness, or
nontrivial-coefficient hypothesis is added.
\end{proposition}
\begin{proof}
\leanok
We give the quotient construction explicitly, using the checked standard
basis rather than replacing the quotient by another ring. More generally,
for a weight $w:V\to M$ into any additive commutative monoid, partition the
basis indices into $J_d=\{\alpha\in J:|\alpha|_w=d\}$ and put
$B_d=\operatorname{span}_R\{b_\alpha:\alpha\in J_d\}$.
The basis-coordinate equivalence identifies $B_d$ with the finitely
supported coefficient functions on $J$ whose support is contained in
$J_d$. Every basis index lies in exactly one fiber. Hence these submodules
span $A$, and a sum from distinct fibers can vanish only when each of its
basis coefficients vanishes. Equivalently, each fiber submodule is disjoint
from the sum of all the others. Thus their sum is internal. Transporting
the supported-coordinate decomposition through the basis equivalence gives
actual direct-sum decomposition data on $A$; each element has only finitely
many nonzero homogeneous components.

These components also agree with the images of the polynomial weighted
homogeneous submodules under $q$. Indeed, a polynomial homogeneous of degree
$d$ is a sum of monomials of weight $d$. Monomials with nonface support
vanish in the quotient, while the others give exactly the displayed basis
vectors of $B_d$. Conversely every such basis vector has its homogeneous
monomial representative. Lemma~\ref{lem:homogeneous} ensures that
$q(\pi_d(p))$ is independent of the chosen polynomial representative $p$:
changing the representative changes its weighted component by an element
of the nonface ideal. These descended projections are exactly the
coordinate-fiber projections just constructed.

For any exponents $\alpha,\beta$, multiplication of monomials gives
$q(x^\alpha)q(x^\beta)=q(x^{\alpha+\beta})$. If the sum exponent has face
support, this is its standard basis vector; otherwise it is zero by ideal
membership. Since weighted degrees add, the product of a basis vector in
$B_d$ and one in $B_e$ lies in $B_{d+e}$ in either case. Bilinearity extends
this to $B_dB_e\subseteq B_{d+e}$. The zero exponent has empty face support,
and its basis vector is $q(1)=1$, of weight zero. Internal decomposition,
the unit condition, and this multiplication law supply the graded-algebra
structure.

For fine weights $w(v)=e_v$, the checked identity
$|\alpha|_w=\alpha$ makes each fiber either a singleton basis index or
empty. This gives precisely the stated $A_\alpha$, including zero for a
nonface exponent. For unit weights, $|\alpha|_w=|\alpha|$ gives $T_n$.
Its basis-index fiber is the union of the fine fibers with
$|\alpha|=n$, so its submodule is their supremum, and their sum is internal
by the same basis independence. Thus the total grading is the asserted
coarsening of the fine grading, with both its direct sum and multiplication
law proved.

For nonnegative integer exponents, $|\alpha|=0$ holds exactly when
$\alpha=0$. Consequently $T_0=A_0=R\,b_0=R\,1$. A scalar multiple of $1$
is the image of that scalar under the algebra map. Proposition~\ref{prop:basis}
already proves this scalar map injective, and the displayed span makes it
surjective onto the degree-zero subalgebra. Restricting the scalar algebra
homomorphism to that subalgebra therefore gives the stated algebra
equivalence, with its value at every scalar explicitly the original scalar
image in $A$.

If $K=\{\varnothing\}$, the only face-supported exponent is zero, so both
gradings are concentrated in degree zero and $A\simeq_R R$. For a ghost
vertex $v$, the fine component at $e_v$ is zero, since its support is the
nonface $\{v\}$; the ambient index set is unchanged. Infinite ambient types
cause no problem because every exponent and every basis expansion has finite
support. The component families, internal gradings, coarsening, zero
components, and scalar algebra equivalence are checked in Lean.
\end{proof}

\begin{proposition}[Ghost vertices]\label{prop:ghost}
\uses{prop:basis,def:dimension,def:homology,def:link}
Passing from the ambient $V$ to $W=\{v:\{v\}\in K\}$ gives an isomorphic
face ring. Links and augmented homology are unchanged after removing ghosts.
For $K=\{\varnothing\}$ the face ring is $k$, even when $V$ is nonempty.
\end{proposition}
\begin{proof}
Every ghost variable is a nonface generator and hence zero in $A$. Inclusion
of the remaining variables and evaluation of ghost variables at zero induce
inverse maps of the two quotients: their composites agree on coefficients
and every variable. The face bases in the chain complexes coincide, as no
face contains a ghost. If there are no actual vertices the only surviving
monomial is $1$.
\end{proof}

\begin{proposition}[Finite type and the augmentation ideal]\label{prop:augmentation}
\uses{def:ring,prop:basis,prop:grading}
\lean{Reisner.Complex.faceRingAugmentation, Reisner.Complex.faceRingAugmentation_mk, Reisner.Complex.faceRingAugmentation_algebraMap, Reisner.Complex.faceRingAugmentation_surjective, Reisner.Complex.faceRingAugmentationIdeal, Reisner.Complex.faceRingPositiveIdeal, Reisner.Complex.idealOfVars_eq_ker_constantCoeff, Reisner.Complex.faceRingAugmentationIdeal_eq_map_idealOfVars, Reisner.Complex.faceRingAugmentation_ker, Reisner.Complex.faceRingResidueEquiv, Reisner.Complex.faceRingResidueEquiv_mk, Reisner.Complex.faceRing_finiteType, Reisner.Complex.faceRing_isNoetherianRing, Reisner.Complex.faceRingAugmentationIdeal_eq_positive, Reisner.Complex.faceRingAugmentationIdeal_isHomogeneous, Reisner.Complex.faceRingAugmentationIdeal_isMaximal, Reisner.Complex.faceRing_variable_eq_zero_of_ghost, Reisner.Complex.faceRingAugmentationIdeal_eq_span_vertices, Reisner.Complex.faceRingAugmentationIdeal_eq_bot_of_no_vertices, Reisner.Complex.faceRingAugmentationEquivOfNoVertices, Reisner.Complex.faceRingAugmentationEquivOfNoVertices_apply}
\leanok
Let $K$ be an augmented complex on any ambient type $V$, let $R$ be a
commutative ring, and put $A=R[K]$, with quotient map
$q:R[x_v:v\in V]\to A$ and total-degree components $T_n$ from
Proposition~\ref{prop:grading}. Evaluation of every variable at zero induces
an augmentation algebra map
\[
 \varepsilon:A\longrightarrow R,\qquad
 \varepsilon(q(p))=[x^0]p.
\]
It satisfies $\varepsilon(r\,1)=r$ and is surjective. Define independently
the variable ideal and the positive-degree ideal by
\[
 \mathfrak m=(q(x_v):v\in V),\qquad
 A_{>0}=\bigl(\,\bigcup_{n>0}T_n\,\bigr).
\]
The parentheses in the second expression mean the ideal generated by the
displayed homogeneous elements. Then
\[
 \ker\varepsilon=\mathfrak m=A_{>0},\qquad
 A/\mathfrak m\simeq_R R.
\]
The residue equivalence sends the class of $a$ to $\varepsilon(a)$.
The ideal $\mathfrak m$ is homogeneous for the total grading; if $R=k$
is a field, it is maximal and the residue algebra is exactly $k$.
If $V$ is finite, $A$ is a finite-type $R$-algebra. If moreover $R$ is
Noetherian, then $A$ is Noetherian. In particular these conclusions hold
for a finite ambient type over any field, and $\mathfrak m$ is precisely
the homogeneous maximal ideal used in Proposition~\ref{prop:dimension}.

For a ghost vertex $v$, $q(x_v)=0$, so $\mathfrak m$ can equivalently be
generated by the classes indexed only by actual vertices. This omits zero
generators inside the same quotient ring. If $K=\{\varnothing\}$, then
$\mathfrak m=0$ and $\varepsilon$ itself is an algebra equivalence
$A\simeq_R R$. No purity, actual-vertex, nonempty-ambient, or no-ghost
hypothesis is imposed; finiteness is needed only for the finite-type and
Noetherian conclusions.
\end{proposition}
\begin{proof}
\leanok
The constant-coefficient map on the polynomial ring is a ring homomorphism
and fixes the coefficient ring. For $p\in I_K$, a nonzero coefficient at
exponent zero would put zero in the polynomial support. The checked
membership theorem would then say that its empty vertex support is a
nonface, contradicting the empty-face witness. Thus the constant-coefficient
map kills $I_K$ and descends through the actual quotient to the stated
algebra map. Its value on scalar classes is the original scalar, proving
surjectivity and compatibility with the checked scalar embedding.

In the polynomial ring, put $J=(x_v:v\in V)$. A polynomial is in $J$
exactly when its constant coefficient is zero: every nonconstant monomial
has a variable factor, while every multiple of a variable has zero constant
coefficient. The existing variable-ideal membership theorem gives this
coefficient characterization directly, without a finite-vertex hypothesis.
Hence $J$ is the kernel of the polynomial constant-coefficient map.
Its image under $q$ is the ideal generated by the variable classes, namely
$\mathfrak m$. Every element of $A$ has a polynomial representative; the
augmentation of that representative vanishes exactly when the representative
is in $J$. This proves $\ker\varepsilon=q(J)=\mathfrak m$.

Each variable class belongs to $T_1$, so the variable ideal is contained
in $A_{>0}$. Conversely, by the polynomial-image comparison in
Proposition~\ref{prop:grading}, every element of $T_n$ has a polynomial
representative homogeneous of total degree $n$. For $n>0$ its exponent-zero
coefficient is zero, so its augmentation is zero. The kernel is an ideal;
it therefore contains the ideal generated by all these positive components.
Together with the kernel calculation this proves
$\mathfrak m=A_{>0}$. These generators are homogeneous, so their ideal is
homogeneous. Surjectivity and the first isomorphism theorem give the stated
residue algebra equivalence, with the displayed value formula. When the
target $R$ of the augmentation is a field, this quotient is a field and
$\mathfrak m$ is maximal.

For finite $V$, the polynomial algebra is generated by its finitely many
variables, and its quotient is a finite-type algebra. If $R$ is Noetherian,
Hilbert's basis theorem and passage to a quotient make $A$ Noetherian.
These are the standard finite-type and quotient results already available
in the pinned library; this step makes no Krull-dimension assertion.

A ghost singleton is a nonface, so its variable class vanishes by checked
monomial membership. Removing these zero elements from the generating set
does not change its ideal. If the empty face is the only face, every
singleton is a ghost, whence $\mathfrak m=0$. The augmentation then has
zero kernel and is both injective and surjective, giving the stated algebra
equivalence. This argument includes an empty ambient type and arbitrarily
many ghost vertices; it neither changes the ambient type nor claims a
ghost-removal comparison for links or homology. The finite-type and
Noetherian statements, augmentation and ideal identifications, homogeneity,
maximality, residue equivalence, and these ghost and empty-face-only
consequences are checked in Lean.
\end{proof}

\begin{proposition}[Coordinate primes and facet intersection]\label{prop:facet-primes}
\uses{def:complex,def:ring,prop:basis}
\lean{Reisner.coordinateIdeal, Reisner.mem_coordinateIdeal_iff, Reisner.coordinateRestriction, Reisner.coordinateRestriction_ker, Reisner.coordinateRestriction_surjective, Reisner.coordinateQuotientEquiv, Reisner.coordinateQuotientEquiv_mk, Reisner.coordinateRestriction_X_inside, Reisner.coordinateRestriction_X_outside, Reisner.coordinateIdeal_isPrime, Reisner.coordinateQuotient_ringKrullDim, Reisner.Complex.facets, Reisner.Complex.facets_finite, Reisner.Complex.exists_facet_superset, Reisner.Complex.facets_nonempty, Reisner.Complex.stanleyReisnerIdeal_eq_iInf_coordinateIdeal, Reisner.Complex.exists_facet_coordinateIdeal_le, Reisner.Complex.ghost_not_mem_facet, Reisner.Complex.facets_eq_singleton_empty_of_no_vertices}
\leanok
Let $V$ be a finite ambient type, $k$ a field, $K$ an augmented complex,
and $S=k[x_v:v\in V]$. Write
\[
 \operatorname{Fac}(K)=
 \{F\in K:\ \forall G\in K,\ F\subseteq G\Longrightarrow G\subseteq F\}.
\]
Thus facets are maximal faces, with the empty face allowed. The facet
type is finite and nonempty, and every face is contained in a facet.
For any finite set $F\subseteq V$, independently of whether it is a face,
define the coordinate ideal by its variable generators:
\[
 P_F=(x_v:v\notin F)\subseteq S.
\]
For every polynomial $p$,
\[
 p\in P_F\quad\Longleftrightarrow\quad
 \forall\alpha\in\operatorname{supp}p,\quad
       \operatorname{supp}\alpha\not\subseteq F.
\]
Equivalently, each nonzero term has a variable outside $F$ in its support.
Setting those outside variables to zero and retaining the variables of $F$
induces the explicit algebra equivalence
\[
 S/P_F\simeq_k k[x_v:v\in F].
\]
The ideal $P_F$ is prime, and $\dim(S/P_F)=|F|$. In Lean the latter is
an equality for ring Krull dimension with the natural-number cardinality
cast into its extended dimension type. The actual nonface-generated ideal
satisfies the equality for arbitrary polynomials
\[
 I_K=\bigcap_{F\in\operatorname{Fac}(K)}P_F.
\]
Every prime ideal $Q$ of $S$ containing $I_K$ contains $P_F$ for at least
one facet $F$.

Coordinate membership holds over any commutative semiring. The coordinate
quotient equivalence holds over any commutative ring, and the coordinate
ideal is prime over any integral domain. These coordinate statements, and
the quotient dimension formula over a field, need no finiteness of the
ambient type because $F$ itself is finite. Finite ambient vertices are
used for facet extension and the finite intersection and prime-containment
arguments. No purity or no-ghost hypothesis is imposed.
\end{proposition}
\begin{proof}
\leanok
The facet predicate is maximality in the inclusion order on finite faces.
There are finitely many subsets of a finite ambient type, so the facet
subtype is finite. Every face has a maximal face above it in this finite
partially ordered set. Applying this to the empty face also shows that
the facet type is nonempty. This uses the augmented convention, including
when the empty face is the only face.

The variable-ideal membership theorem says that $p\in P_F$ exactly when
each exponent in its polynomial support has a nonzero coordinate at some
$v\notin F$. This is precisely the displayed support condition.
Define $r_F:S\to k[x_v:v\in F]$ by sending $x_v$ to its corresponding
variable if $v\in F$, and to zero otherwise, while fixing coefficients.
The polynomial inclusion from the variable subtype $F$ back to $V$ is a
right inverse to $r_F$, so $r_F$ is surjective.

Under $r_F$, the terms whose exponent support lies in $F$ retain their
coefficients, and all other terms vanish. Distinct exponents supported in
$F$ remain distinct after this restriction, so there is no cancellation
between different surviving terms. Therefore $r_F(p)=0$ holds exactly when
every exponent in the support has a vertex outside $F$. The membership
criterion proves $\ker r_F=P_F$. The first isomorphism theorem gives the
stated equivalence of the actual ideal quotient, sending a quotient class
to its specialized polynomial. Its target is a polynomial ring over a
domain, hence a domain; equivalently, the kernel of $r_F$ is prime. Over
a field the finite-variable polynomial dimension theorem gives
$\dim k[x_v:v\in F]=\dim k+|F|=|F|$, and the quotient equivalence transfers
this dimension to $S/P_F$.

We now compare arbitrary polynomial terms. If $p\in I_K$, every exponent
in its support has nonface vertex support by Proposition~\ref{prop:basis}.
Such a support cannot lie in a facet, since downward closure would make
it a face. The coordinate membership criterion therefore puts $p$ in
every $P_F$. Conversely, if $p$ lies in every facet coordinate ideal and
a term of $p$ has face support $T$, extend $T$ to a facet $F$. That term's
support is contained in $F$, contradicting membership in $P_F$. Thus all
terms have nonface support and $p\in I_K$. This proves the intersection
equality for polynomials, including possible cancellation among other terms.

The product of the finitely many ideals $P_F$ is contained in their
intersection $I_K$. If a prime $Q$ contains $I_K$, it contains this product.
Primality for a finite ideal product then gives a factor $P_F\subseteq Q$.
The pinned finite-intersection lemma is exactly this product-and-primality
argument, so it can be reused directly.

A ghost vertex is in no face, hence in no facet; its variable belongs to
every $P_F$. If all singleton faces are absent, downward closure forces
the empty face to be the only face. It is then the unique facet. The
intersection reduces to $I_K=P_{\varnothing}$, the ideal of all variables;
the coordinate quotient has no variables and is $k$, of dimension zero.
These are statements about facets of the given augmented complex, not an
assumed purity condition or a change of ambient vertices. The coordinate
ideals and their quotient dimensions, facet extension and intersection,
prime containment, and these facet conventions are checked in Lean.
\end{proof}

\begin{proposition}[Dimension]\label{prop:dimension}
\uses{prop:facet-primes,prop:augmentation,prop:basis,prop:grading,def:dimension}
\lean{Reisner.X_mem_coordinateIdeal_iff, Reisner.coordinateIdeal_le_iff, Reisner.coordinateIdeal_empty, Reisner.exists_coordinatePrimeChain, Reisner.Complex.stanleyReisnerIdeal_le_coordinateIdeal, Reisner.Complex.faceRing_ringKrullDim_le, Reisner.Complex.exists_faceRingPrimeChain, Reisner.Complex.card_le_faceRingAugmentationIdeal_height, Reisner.Complex.maxFaceCard_le_faceRingAugmentationIdeal_height, Reisner.Complex.faceRingAugmentationIdeal_isPrime, Reisner.Complex.faceRing_ringKrullDim, Reisner.Complex.faceRingAugmentationIdeal_height, Reisner.Complex.faceRing_atAugmentation_ringKrullDim}
\leanok
The ring $A$ is a finitely generated Noetherian $k$-algebra and
$\dim A=\dim A_{\mathfrak m}=d(K)$, where
$\mathfrak m=(\bar x_v:v\in V)=A_{>0}$ is a maximal ideal and
$A/\mathfrak m=k$.
\end{proposition}
\begin{proof}
\leanok
Polynomial rings over a field are Noetherian; a quotient is Noetherian and
finitely generated. Evaluation at zero gives the stated residue field.
These assertions, maximality of $\mathfrak m$, and its identification with
the positive-degree ideal are already checked in
Proposition~\ref{prop:augmentation}.

For the upper bound, take any finite strict chain of primes containing
$I_K$. By Proposition~\ref{prop:facet-primes}, its first prime contains a
facet coordinate ideal $P_F$, hence every prime in the chain contains $P_F$.
The quotient prime-spectrum order isomorphism views the same chain, with
the same length, in $\operatorname{Spec}(S/P_F)$. The checked coordinate
quotient formula bounds that length by $|F|\leq d(K)$. Taking the supremum
over all chains gives $\dim A\leq d(K)$.

For the lower bound, first observe that $x_v\in P_G$ exactly when
$v\notin G$. Thus $P_G\subseteq P_H$ exactly when $H\subseteq G$.
For every finite $F$, removing its vertices one at a time gives a strict
coordinate-prime chain from $P_F$ to $P_{\varnothing}=(x_v:v\in V)$ of
length $|F|$. More explicitly, the empty set gives a chain of length zero;
if $v\notin G$, prepend $P_{G\cup\{v\}}$ to the chain for $G$.
The new inclusion is strict, witnessed by $x_v$, and increases the length
by one. If $F$ is a face, termwise membership and downward closure show
$I_K\subseteq P_F$, so the entire chain descends to $A$ through the
quotient prime-spectrum order isomorphism without losing strictness.
Its final prime is $\mathfrak m$, since $P_{\varnothing}$ maps to the
augmentation ideal. This proves $|F|\leq\operatorname{ht}\mathfrak m$
for every face, even without finite ambient $V$.

A face, and indeed a facet, attains the finite maximum $d(K)$. Consequently
\[
 d(K)\leq\operatorname{ht}\mathfrak m\leq\dim A\leq d(K).
\]
The standard localization prime correspondence identifies
$\dim A_{\mathfrak m}$ with $\operatorname{ht}\mathfrak m$, proving both
required equalities. Lean uses the actual localization
\texttt{Localization.AtPrime} and the pinned localization-height theorem;
the ring dimensions lie in \texttt{WithBot} $\mathbb N\cup\{\infty\}$,
with $d(K)$ cast from $\mathbb N$. No purity or no-ghost hypothesis occurs.
For $K=\{\varnothing\}$ the chain has length zero, both ring dimensions
are zero, and the simplicial dimension remains $-1$.
The foundational algebra used here is Hilbert's basis theorem,
quotient/localization prime correspondence, and the dimension of a polynomial
ring over a field; see Stacks, Sections 10.31, 10.60 and 10.114
\cite{stacks-dimension}.
\end{proof}

\section{Local depth and local cohomology}

\begin{lemma}[The unique maximal homogeneous ideal of a face ring]\label{lem:face-ring-graded-maximal}
\uses{prop:grading,prop:augmentation}
Let $K$ be a simplicial complex. For the total $\mathbb N$-grading on its
face ring over any commutative coefficient ring, the irrelevant ideal is
the augmentation ideal. Over a field $k$, write $A=k[K]$ and
$\mathfrak m=A_+$. Every proper homogeneous ideal is contained in
$\mathfrak m$. In particular $\mathfrak m$ is the unique
maximal element among the proper homogeneous ideals.
\lean{Reisner.Complex.faceRing_irrelevant_eq_augmentation, Reisner.Complex.faceRing_homogeneousIdeal_le_augmentation, Reisner.Complex.faceRing_isCoatom_homogeneousIdeal_iff}
\leanok
\end{lemma}
\begin{proof}
\leanok
Both $A_+$ and $\mathfrak m$ are generated by the positive-degree
components. An element lies in $A_+$ precisely when its degree-zero
component vanishes. If $x$ lies in a homogeneous ideal $I$, its
degree-zero component also belongs to $I$. Since $A_0=k$, a nonzero
degree-zero component would give a unit in $I$. Thus every element
of a proper homogeneous ideal belongs to $A_+$. The augmentation
quotient is the field $k$, so $\mathfrak m$ is itself proper, and it
therefore is the greatest proper homogeneous ideal.
\end{proof}

\begin{theorem}[Graded local test: external foundation]\label{found:graded-cm}
\uses{def:cm,found:cm-localization}
For a finitely generated positively graded commutative $k$-algebra $B$ with
$B_0=k$, put $\mathfrak n=B_{>0}$. Then $B$ is Cohen--Macaulay if and only
if $B_{\mathfrak n}$ is Cohen--Macaulay.
\end{theorem}
Bruns--Herzog, item 2.1.27 \cite{bh}, tests the Cohen--Macaulay property
of a Noetherian graded ring at all its maximal homogeneous ideals, as the
authors confirm in their corrections, page 1, note to 2.1.29
\cite{bh-errata}. In the displayed specialization finite generation over
$k$ gives Noetherianity, and the positive grading with $B_0=k$ gives
a single maximal homogeneous ideal. For a face ring with finite ambient
vertex set these assumptions follow from Propositions~\ref{prop:grading}
and~\ref{prop:augmentation}, and Lemma~\ref{lem:face-ring-graded-maximal}
identifies the ideal. Thus the converse from the augmentation localization
to all prime localizations is recorded as this external graded-local test,
not as a checked Lean theorem. One route uses the homogeneous core
$\mathfrak p^*$ of each prime: it is a homogeneous prime contained in both
$\mathfrak p$ and $\mathfrak n$, and the ordinary localization theorem for
Cohen--Macaulay local rings transfers the property from $B_{\mathfrak n}$
through $B_{\mathfrak p^*}$ to $B_{\mathfrak p}$. The localization theorem
is recorded separately below.

\begin{theorem}[Cohen--Macaulayness localizes: external foundation]\label{found:cm-localization}
\uses{def:cm}
If $(R,\mathfrak m)$ is a Noetherian Cohen--Macaulay local ring and
$\mathfrak p\subseteq\mathfrak m$ is prime, then $R_{\mathfrak p}$ is a
Cohen--Macaulay local ring.
\end{theorem}
This is Stacks Project, Lemma 10.104.5, tag 00NB
\cite{stacks-cm-localize}. It is not implemented in the pinned Mathlib
and is not treated as a Lean declaration.

\begin{definition}[Local cohomology]\label{def:local-cohomology}
For a commutative ring $B$, an ideal $J$, a $B$-module $M$, and $q\ge0$,
use the standard Ext-colimit model
\[
 H^q_J(M)=\mathop{\mathrm{colim}}_{n\in\mathbb N}
           \operatorname{Ext}^q_B(B/J^n,M).
\]
For Lean, take $\operatorname{Ext}$ from the derived-category Ext API used
by Mathlib's depth development and package it with
\texttt{CategoryTheory.Abelian.linearExtFunctor}. This avoids a conversion
between Mathlib's two Ext implementations: the older left-derived-functor
Ext used by \texttt{Mathlib.Algebra.Homology.LocalCohomology}, and the
derived-category Ext used for depth. Both have the standard Ext meaning, but
no Lean isomorphism between their colimits is claimed here. The Ext-colimit
is the standard local cohomology model; homogeneous grading used later is
transported through the Cech comparison.
\lean{Reisner.depthExtLocalCohomologyDiagram, Reisner.depthExtLocalCohomology}
\leanok
\end{definition}

\begin{lemma}[Vanishing in positive-power Ext stages]\label{lem:depth-ext-stage-vanishing}
\uses{def:local-cohomology}
Let $(B,\mathfrak n)$ be a Noetherian local ring, let $q\ge0$, and let
$r\ge1$. If $q<\operatorname{depth}B$, then
\[
 \operatorname{Ext}^q_B(B/\mathfrak n^r,B)=0.
\]
\lean{Reisner.depthExtLocalCohomology_stage_subsingleton}
\leanok
\end{lemma}

\begin{lemma}[Nonvanishing at the depth stage]\label{lem:depth-ext-stage-nontrivial}
\uses{def:local-cohomology}
Let $(B,\mathfrak n)$ be a Noetherian local ring and let
$q=\operatorname{depth}B$. For each $r\ge1$,
\[
 \operatorname{Ext}^q_B(B/\mathfrak n^r,B)\ne0.
\]
\lean{Reisner.depthExtLocalCohomology_stage_nontrivial}
\leanok
\end{lemma}

\begin{lemma}[Local cohomology vanishes below depth]\label{lem:lc-vanishes-below-depth}
\uses{def:local-cohomology,lem:depth-ext-stage-vanishing}
Let $(B,\mathfrak n)$ be a Noetherian local ring. If
$q<\operatorname{depth}B$, then
\[
 H^q_{\mathfrak n}(B)=0.
\]
\lean{Reisner.depthExtLocalCohomology_vanishes_below_depth}
\leanok
\end{lemma}

\begin{theorem}[Local depth cutoff detected by Ext-colimit local cohomology]\label{thm:depth-cutoff}
\uses{def:local-cohomology,lem:lc-vanishes-below-depth,
      lem:depth-ext-stage-nontrivial}
Let $(B,\mathfrak n)$ be a Noetherian local ring. For every finite cutoff
$d\in\mathbb N$,
\[
 d\le\operatorname{depth}B
 \quad\Longleftrightarrow\quad
 H^q_{\mathfrak n}(B)=0\text{ for every }q<d.
\]
Here local cohomology is the Ext-colimit model defined above. In Lean, zero
is stated as `Subsingleton` of the corresponding module.
\lean{Reisner.le_depth_iff_localCohomology_vanishes}
\leanok
\end{theorem}
\begin{proof}
If $d\le\operatorname{depth}B$, each $q<d$ lies below depth, so every
positive-power Ext stage vanishes and hence so does the colimit. Conversely,
if depth is below $d$, it is a finite natural number $r<d$. The stage
$\operatorname{Ext}^r_B(B/\mathfrak n,B)$ is nonzero. The short exact
sequences for adjacent powers show that every transition in degree $r$ is
injective: for $r>0$, the kernel $\mathfrak n^n/\mathfrak n^{n+1}$ is finite
and supported on $V(\mathfrak n)$, so its Ext groups below $r$ vanish by the
Rees theorem; for $r=0$, use injectivity of precomposition along the quotient
epimorphism. Thus the nonzero class survives in the filtered colimit.
\end{proof}

\begin{lemma}[Principal localization kernel is power torsion]\label{lem:principal-localization-torsion}
Let $R$ be a commutative ring, $M$ an $R$-module, and $f\in R$. The
kernel of the localization map from $M$ to its localization at
$1,f,f^2,\ldots$ consists precisely of the elements $m$ for which
$f^n m=0$ for some $n\ge0$.
\lean{Reisner.principalLocalization_kernel_iff_torsion}
\leanok
\end{lemma}

\begin{lemma}[Degree-zero derived Ext is Hom]\label{lem:ext-zero-quotient-hom}
For an ideal $J$ of a commutative ring $R$ and an $R$-module $M$,
\[
 \operatorname{Ext}^0_R(R/J,M)\cong\operatorname{Hom}_R(R/J,M).
\]
The isomorphism here uses the derived-category Ext model shared with the
depth development.
\lean{Reisner.depthExtLocalCohomology_extZero_quotientHomEquiv}
\leanok
\end{lemma}

\begin{lemma}[Maps from a quotient are annihilated elements]\label{lem:quotient-hom-annihilator}
For an ideal $J$ of a commutative ring $R$ and an $R$-module $M$,
\[
 \operatorname{Hom}_R(R/J,M)\cong
 \{m\in M: a m=0\text{ for every }a\in J\}.
\]
The map evaluates a homomorphism at $1+J$.
\lean{Reisner.depthExtQuotientHomTorsionEquiv}
\leanok
\end{lemma}

\begin{lemma}[Naturality of the degree-zero Ext/Hom identification]\label{lem:ext-zero-naturality}
\uses{lem:ext-zero-quotient-hom}
For a map $f:X\to Y$ and $x\in\operatorname{Ext}^0_R(Y,M)$, the degree-zero
Ext/Hom equivalences send precomposition by $f$ to composition by $f$:
\[
  e_X(f^*x)=f\circ e_Y(x).
\]
This checked compatibility identifies the transition maps on the quotient
stages with the corresponding inclusions of annihilator submodules.
\lean{Reisner.depthExt_zero_precomp_naturality}
\leanok
\end{lemma}

\begin{lemma}[Degree-zero transitions are annihilator inclusions]\label{lem:ext-zero-torsion-transition}
\uses{lem:ext-zero-quotient-hom,lem:quotient-hom-annihilator,
      lem:ext-zero-naturality}
For every ideal $I$, module $M$, and $n\ge0$, the transition from
$\operatorname{Ext}^0_R(R/I^n,M)$ to $\operatorname{Ext}^0_R(R/I^{n+1},M)$
corresponds under Ext$^0\cong\operatorname{Hom}$ and evaluation at $1$ to the
inclusion of the submodule annihilated by $I^n$ into the one annihilated by
$I^{n+1}$. For $I=(f)$, stage $n$ is precisely the elements killed by $f^n$.
\lean{Reisner.depthExtLocalCohomology_transition_torsion_natural}
\lean{Reisner.principalPowerTorsion_eq}
\leanok
\end{lemma}
\begin{lemma}[Degree-zero Ext-colimit is a power-torsion supremum]\label{lem:ext-zero-power-torsion-sup}
\uses{def:local-cohomology,lem:ext-zero-torsion-transition}
For every ideal $I$ of a commutative ring $R$ and every $R$-module $M$, the
Ext-colimit degree-zero local cohomology module is linearly equivalent to
\[
  \sup_{n\in\mathbb N}\operatorname{Ann}_M(I^n).
\]
The supremum is the union of the nested annihilator submodules, and stage zero
is included.
\lean{Reisner.depthExtLocalCohomology_zero_torsionSupEquiv}
\leanok
\end{lemma}

\begin{theorem}[Principal degree-zero local cohomology is the localization kernel]\label{thm:principal-degree-zero-localization}
\uses{lem:ext-zero-power-torsion-sup,lem:principal-localization-torsion}
Let $R$ be a commutative ring, $M$ an $R$-module, and $f\in R$. For the
Ext-colimit model used in the depth development, there is a linear equivalence
\[
 H^0_{(f)}(M)\cong\ker\bigl(M\longrightarrow M_f\bigr),
\]
where localization is at $1,f,f^2,\ldots$. No nonzero or regularity hypothesis
on $f$ is required; the exponent-zero stage and $f=0$ are included.
\lean{Reisner.principalExtZeroKernelEquiv}
\leanok
\end{theorem}
\begin{proof}
At stage $n$, Ext$^0_R(R/(f^n),M)$ identifies with maps from $R/(f^n)$,
then with elements killed by $f^n$. The checked transition compatibility
shows the colimit is the supremum of these nested submodules. A class belongs
to that supremum exactly when it is killed by some power of $f$, which is
precisely the localization kernel criterion.
\end{proof}

\begin{lemma}[Finite powers annihilating generators give ideal-power torsion]\label{lem:finite-generator-power-torsion}
\uses{lem:principal-localization-torsion}
Let $s$ be a finite set of indices and let $I=(f_i : i\in s)$. If for each
$i\in s$ some power $f_i^{k_i}$ annihilates $m\in M$, then
$I^{1+\sum_{i\in s} k_i}m=0$. Consequently, for arbitrary $R$ and $M$,
\[
  \bigcup_{n\ge 0}\operatorname{Ann}_M(I^n)
  =\bigcap_{i\in s}\ker(M\to M_{f_i}).
\]
The empty index set is allowed: $I=0$, and the stage-zero annihilator is all
of $M$, matching the empty intersection.
\lean{Reisner.finiteGeneratorsLocalizationKernel_eq_powerTorsionSup}
\leanok
\end{lemma}
\begin{proof}
The forward inclusion uses $f_i^n\in I^n$. For the reverse inclusion, choose
$k_i$ from the principal localization kernel criterion. Induct on the finite
generator set, using $(J+K)^{a+b}\subseteq J^a+K^b$ and
$(f_i)^{k_i}\subseteq\operatorname{Ann}(m)$. The additional one in the total
exponent handles the empty set and keeps exponent zero valid for zero
generators.
\end{proof}

\begin{theorem}[Finite-generator degree-zero localization comparison]\label{thm:finite-generator-degree-zero-localization}
\uses{lem:ext-zero-power-torsion-sup,lem:finite-generator-power-torsion}
For $I=(f_i:i\in s)$, the degree-zero Ext-colimit local cohomology module
used in this development is linearly equivalent to
\[
  \bigcap_{i\in s}\ker(M\to M_{f_i}).
\]
This is a comparison with the simultaneous localization kernel; no
higher-degree Cech comparison is asserted here.
\lean{Reisner.finiteGeneratorsExtZeroLocalizationKernelEquiv}
\leanok
\end{theorem}
\begin{proof}
Compose the checked equivalence from the Ext-colimit to the supremum of the
annihilator submodules with the finite-generator kernel identity.
\end{proof}

\begin{lemma}[Degree-zero Čech cycles are simultaneous localization kernels]\label{lem:cech-degree-zero-kernel}
\uses{def:complex,lem:localization-kernel}
Let $K$ be a simplicial complex on a finite linearly ordered vertex set $V$,
and let $R$ be any commutative ring. Under the canonical identification of
the degree-zero Čech cochains with the face ring $R[K]$, the kernel of the
outgoing degree-zero differential is
\[
  \bigcap_{v\in V}\ker\bigl(R[K]\longrightarrow R[K]_{\bar x_v}\bigr).
\]
The intersection includes ghost vertices and singleton nonfaces, whose
localizations are zero rings.
\lean{Reisner.Complex.cechZeroCyclesLocalizationEquiv}
\leanok
\end{lemma}
\begin{proof}
The degree-zero term has the unique empty-face coordinate. At a singleton
target coordinate, the alternating sum has one term with sign $+1$, namely
the localization map. The degree-one index type consists exactly of these
singletons, so the differential vanishes precisely when every localization
map vanishes. The argument includes ghost vertices, where the target ring
is subsingleton.
\end{proof}

\begin{theorem}[Degree-zero Ext-colimit local cohomology is Čech cycles]\label{thm:ext-zero-cech-cycles}
\uses{thm:finite-generator-degree-zero-localization,lem:cech-degree-zero-kernel}
Let $K$ be a simplicial complex on a finite ordered vertex set $V$ and let
$R$ be a commutative ring. For the face ring $A=R[K]$ and its augmentation
ideal $\mathfrak m=(\bar x_v:v\in V)$, restrict scalars on the degree-zero
Ext-colimit module along $R\to A$. Then
\[
  \mathrm{colim}_{n\in\mathbb N}\operatorname{Ext}^0_A(A/\mathfrak m^n,A)
  \cong \ker\bigl(C^0(K;A)\longrightarrow C^1(K;A)\bigr)
\]
as $R$-modules.
\lean{Reisner.Complex.extZeroCechCyclesEquiv}
\leanok
\end{theorem}
\begin{proof}
Apply the finite-generator degree-zero localization equivalence to the
singleton face-ring generators. Their span is the augmentation ideal. The
degree-zero Čech-kernel equivalence identifies the resulting simultaneous
localization kernel with the kernel of the outgoing Čech differential.
Restriction of scalars makes both maps $R$-linear.
\end{proof}

\begin{theorem}[Degree-zero Ext-colimit and Čech cohomology]\label{thm:ext-zero-cech-cohomology}
\uses{thm:ext-zero-cech-cycles,lem:cech-complex}
Let $K$ be a simplicial complex on a finite linearly ordered set $V$ and
let $R$ be any commutative ring. Write $A=R[K]$,
$\mathfrak m=(\bar x_v:v\in V)$, and $C^\bullet$ for the extended Čech
complex on these generators, including the ghost variables. There are
$R$-linear isomorphisms
\[
  \mathrm{colim}_{n\in\mathbb N}\operatorname{Ext}^0_A(A/\mathfrak m^n,A)
  \ \cong\ \ker(d^0:C^0\longrightarrow C^1)
  \ \cong\ H^0(C^\bullet),
\]
where the Ext-colimit carries the scalar action restricted along $R\to A$.
The second map sends a cycle to its cohomology class.
\lean{Reisner.Complex.cechZeroCyclesHomologyEquiv, Reisner.Complex.extZeroCechCohomologyEquiv}
\leanok
\end{theorem}
\begin{proof}
Theorem~\ref{thm:ext-zero-cech-cycles} gives the first isomorphism.
The module kernel of $d^0$ and the cycles object of $C^\bullet$ are
kernels of the same map $C^0\to C^1$. Their universal properties give a
canonical isomorphism respecting the inclusions into $C^0$.
The natural-number cochain shape allows a differential $C^i\to C^j$
only when $j=i+1$, so every differential into degree zero vanishes.
Thus the boundary submodule is zero and the canonical projection from
cycles to the cohomology object is an isomorphism.
\end{proof}

\begin{lemma}[Cofinal powers of face-ring generators]\label{lem:generator-power-ideals}
\uses{prop:augmentation}
Let $K$ be a simplicial complex on a finite vertex set $V$, let $R$ be any
commutative ring, and put $A=R[K]$ and $\mathfrak m=(\bar x_v:v\in V)$.
For each $n\in\mathbb N$, set $J_n=(\bar x_v^n:v\in V)$. Then
\[
  J_n\subseteq\mathfrak m^n,\qquad
  \mathfrak m^{|V|n+1}\subseteq J_n,\qquad J_1=\mathfrak m.
\]
Both decreasing ideal diagrams, $n\mapsto\mathfrak m^n$ and
$n\mapsto J_n$, are initial in the category of ideals $L$ with
$\mathfrak m\subseteq\sqrt L$. No Noetherian or nontriviality
hypothesis on $R$ is required. When $V=\varnothing$, every $J_n=0$,
including $J_0$, while $\mathfrak m^0=A$.
\lean{Reisner.Complex.faceRingGeneratorPowerIdeal, Reisner.Complex.faceRingGeneratorPowerIdeal_one, Reisner.Complex.faceRingGeneratorPowerIdeal_le_pow, Reisner.Complex.faceRingAugmentationIdeal_pow_le_generatorPowerIdeal, Reisner.Complex.faceRingGeneratorPowerIdeal_antitone, Reisner.Complex.faceRingAugmentationIdeal_fg, Reisner.Complex.faceRingAugmentationIdeal_le_generatorPowerIdeal_radical, Reisner.Complex.faceRingAugmentationPowers_initial, Reisner.Complex.faceRingGeneratorPowers_initial}
\leanok
\end{lemma}
\begin{proof}
Each generator $\bar x_v^n$ lies in $\mathfrak m^n$. For the other
bound, apply the finite-generator power estimate to the family
$(\bar x_v)$ with every exponent equal to $n$. That estimate follows
by induction on the generator family from
$(I+J)^{a+b}\subseteq I^a+J^b$; the extra one handles the empty family.
The augmentation ideal is finitely generated. Thus if
$\mathfrak m\subseteq\sqrt L$, some power $\mathfrak m^r$ lies in $L$,
and then $J_r\subseteq L$ as well. The corresponding costructured-arrow
categories are nonempty, and any two objects are connected by the arrow
from the larger exponent to the smaller. This proves initiality for
both diagrams, including zero generators and the empty vertex set.
\end{proof}

\begin{theorem}[Generator-power stages for derived Ext local cohomology]\label{thm:generator-power-ext-colimit}
\uses{def:local-cohomology,lem:generator-power-ideals}
With $A$, $\mathfrak m$ and $J_n$ as above, for every cohomological degree
$q\in\mathbb N$ there is an $A$-linear isomorphism
\[
 \operatorname*{colim}_{n}\operatorname{Ext}^q_A(A/\mathfrak m^n,A)
 \ \cong\
 \operatorname*{colim}_{n}\operatorname{Ext}^q_A(A/J_n,A).
\]
Both sides use the derived-category Ext implementation of the project.
The maps come from quotient maps as the exponent increases. The power
index $n$ and the cohomological degree $q$ are independent; no shift
in degree is made. This is a normalization of the stage system and
does not assert a positive-degree comparison with Čech cohomology.
\lean{Reisner.Complex.faceRingGeneratorPowersDiagram, Reisner.Complex.faceRingGeneratorPowerExtDiagram, Reisner.Complex.faceRingLocalCohomologyIsoGeneratorPowerExtColimit}
\leanok
\end{theorem}
\begin{proof}
Form the derived Ext diagram with coefficient module $A$ on the opposite
of the category of ideals $L$ with $\mathfrak m\subseteq\sqrt L$.
The two opposite initial functors of
Lemma~\ref{lem:generator-power-ideals} are final, so their restrictions
have canonically isomorphic colimits. These restricted diagrams are
definitionally the two displayed quotient Ext diagrams. Finally,
evaluation at $A$ commutes with the functor-category colimit defining
the project's local-cohomology functor.
\end{proof}

\begin{lemma}[Derived Ext vanishing from a scalar]\label{lem:localized-ext-scalar}
Let $A$ be a commutative ring and $M,N$ $A$-modules. If some $r\in A$
acts by zero on $M$ and acts invertibly on $N$, then
\[
  \operatorname{Ext}^q_A(M,N)=0\qquad(q\in\mathbb N).
\]
In particular, if $r\in J$ and its image is a unit in a commutative
$A$-algebra $B$, then $\operatorname{Ext}^q_A(A/J,B)=0$ for every $q$.
\lean{Reisner.derivedExt_eq_zero_of_smul_id_eq_zero, Reisner.derivedExt_eq_zero_of_annihilator_isUnit, Reisner.derivedExt_quotient_eq_zero_of_mem_of_isUnit}
\leanok
\end{lemma}
\begin{proof}
The derived Ext module structure identifies multiplication by $r$ induced
from the source with multiplication by $r$ induced from the target. The
source map is zero, while the target map is an isomorphism. Its induced Ext
map is therefore both zero and injective, so every Ext class vanishes. For
$A/J$, membership $r\in J$ says exactly that $r$ annihilates the quotient.
\end{proof}

\begin{theorem}[Ext-colimit vanishing on nonempty Čech rows]\label{thm:localized-ext-rows}
\uses{lem:localized-ext-scalar,def:local-cohomology}
Let $K$ be a simplicial complex on an ambient vertex type $V$, let $R$ be
any commutative ring, put $A=R[K]$ and $\mathfrak m=(\bar x_v:v\in V)$.
For every nonempty finite index $F\subseteq V$, exponent $n\in\mathbb N$
and cohomological degree $q\in\mathbb N$,
\[
  \operatorname{Ext}^q_A(A/\mathfrak m^n,A_{x_F})=0,
  \qquad H^q_{\mathfrak m}(A_{x_F})=0,
\]
where $A_{x_F}$ inverts the product of the variables indexed by $F$ and
local cohomology is the project's derived Ext-colimit model. This includes
$n=0$, ghost variables, and nonface indices; no Noetherian or nontriviality
hypothesis is required. The empty index $F=\varnothing$ is intentionally
excluded.
\lean{Reisner.Complex.faceRingLocalization_variable_isUnit, Reisner.Complex.faceRingLocalization_augmentationPower_ext_eq_zero, Reisner.Complex.faceRingLocalization_generatorPower_ext_eq_zero, Reisner.Complex.faceRingLocalization_augmentationPower_ext_isZero, Reisner.Complex.faceRingLocalization_localCohomology_isZero}
\leanok
\end{theorem}
\begin{proof}
Choose $v\in F$. Its class belongs to $\mathfrak m$, so $\bar x_v^n$
annihilates $A/\mathfrak m^n$. Since $\bar x_v$ divides the inverted
product $x_F$, it is a unit in $A_{x_F}$. Lemma~\ref{lem:localized-ext-scalar}
gives the stage vanishing for every $q$, including $n=0$. The result is
valid even when $F$ is a nonface: then the localization is the zero ring,
where the same unit argument applies. If every stage is zero, the colimit
is zero, proving the local-cohomology assertion.
\end{proof}

\begin{theorem}[Acyclicity of positive Čech terms]\label{thm:cech-acyclic-terms}
\uses{def:local-cohomology,thm:localized-ext-rows}
Let $K$ be a simplicial complex on a finite vertex set $V$, let $R$ be any
commutative ring, and let $A=R[K]$ with augmentation ideal $\mathfrak m$.
The functor $H^q_{\mathfrak m}(-)$ defined by the project's derived
Ext-colimit is additive and preserves finite products. Consequently, for
every positive Čech degree $p>0$ and every $q\in\mathbb N$,
\[
  \operatorname{Ext}^q_A(A/\mathfrak m^n,C^p(K;A))=0\quad(n\in\mathbb N),
  H^q_{\mathfrak m}(C^p(K;A))=0,
  \qquad C^p(K;A)=\prod_{|F|=p} A_{x_F}.
\]
No Noetherian hypothesis is needed for this acyclicity statement.
\lean{Reisner.depthExtTargetFunctor_additive, Reisner.depthExtTargetFunctor_preservesFiniteProducts, Reisner.depthExtLocalCohomology_additive, Reisner.depthExtLocalCohomology_preservesFiniteProducts, Reisner.Complex.faceRingCechCochain_augmentationPower_ext_isZero, Reisner.Complex.faceRingCechCochain_localCohomology_isZero}
\leanok
\end{theorem}
\begin{proof}
The target Ext functor and its colimit are additive, hence preserve finite
products. In positive Čech degree every index $F$ is nonempty.
Theorem~\ref{thm:localized-ext-rows} makes each localized factor Ext-acyclic
at every stage, so finite-product preservation gives the stagewise vanishing
for $C^p$. Taking the colimit gives the stated local-cohomology vanishing.
\end{proof}

\begin{theorem}[The fixed-stage Ext-applied Čech complex]\label{thm:cech-ext-complex}
\uses{def:local-cohomology,thm:cech-acyclic-terms}
Let $K$ be a finite simplicial complex, $R$ a commutative ring,
$A=R[K]$, and $\mathfrak m=(\bar x_v:v\in V)$. For each fixed $n,q$,
first apply $\operatorname{Ext}^q_A(A/\mathfrak m^n,-)$ termwise to the
face-ring-linear extended Čech complex of $A$. The resulting complex has
zero terms in every positive Čech degree, and hence zero homology there.
Its degree-zero homology is canonically $\operatorname{Ext}^q_A(A/\mathfrak m^n,A)$.
This is a fixed-stage complex statement; it does not assert that termwise
application of $\operatorname{Ext}^q$ commutes with Čech cohomology. The
comparison with the colimit and the double-complex identification are still
outstanding. Separately, for the cofinal generator-power ideals
$J_n=(\bar x_v^n:v\in V)$, the same positive-term and positive-homology
vanishing holds for $\operatorname{Ext}^q_A(A/J_n,-)$. These are the stages
of the checked local-cohomology Ext colimit. Their quotient-induced transition
maps are now assembled into a coherent diagram of complexes. The degree-zero
term diagram is naturally isomorphic to the established generator-power Ext
diagram; after taking the colimit, its degree-zero term is the local-cohomology
module. Every positive term, and hence positive homology, of the colimit
complex vanishes. This colimit complex is obtained by applying each fixed
$\operatorname{Ext}^q$ functor to the Čech cochains before taking the stage
colimit. It is not the original Čech complex and does not identify its
positive-degree cohomology with local cohomology. The double-complex or
acyclic-resolution comparison remains open.
\lean{Reisner.Complex.faceRingLocalizationMapA, Reisner.Complex.faceRingLocalizationMapA_algebraMap, Reisner.Complex.faceRingLocalizationMapA_toRingHom, Reisner.Complex.cechDifferentialA, Reisner.Complex.cechDifferentialA_apply, Reisner.Complex.cechDifferentialA_comp_eq_zero, Reisner.Complex.faceRingCechComplex, Reisner.Complex.faceRingCechExtComplex, Reisner.Complex.faceRingCechExtComplex_X_isZero, Reisner.Complex.faceRingCechExtComplex_exactAt, Reisner.Complex.faceRingCechExtComplex_homology_isZero, Reisner.Complex.faceRingLocalizationEmptyEquivA, Reisner.Complex.cechZeroFaceRingEquivA, Reisner.Complex.faceRingCechExtZeroTermIso, Reisner.Complex.faceRingCechExtComplex_zeroTermIso, Reisner.Complex.faceRingCechExtComplex_zeroHomologyIso, Reisner.Complex.faceRingCechExtComplex_zeroExtIso, Reisner.Complex.faceRingCechCochain_generatorPower_ext_isZero, Reisner.Complex.faceRingGeneratorPowerCechExtComplex, Reisner.Complex.faceRingGeneratorPowerCechExtComplex_X_isZero, Reisner.Complex.faceRingGeneratorPowerCechExtComplex_homology_isZero, Reisner.Complex.faceRingGeneratorPowerExtFunctorDiagram, Reisner.Complex.faceRingGeneratorPowerCechExtDiagram, Reisner.Complex.faceRingGeneratorPowerCechExtDiagram_obj, Reisner.Complex.faceRingGeneratorPowerCechExtDiagram_map_f, Reisner.Complex.faceRingGeneratorPowerCechExtDiagram_evalIso, Reisner.Complex.faceRingGeneratorPowerCechExtDiagram_zeroTermIso, Reisner.Complex.faceRingGeneratorPowerCechExtDiagram_X_isZero, Reisner.Complex.faceRingGeneratorPowerCechExtDiagram_zeroTermColimitIso, Reisner.Complex.faceRingGeneratorPowerCechExtDiagram_colimitZeroTermIso, Reisner.Complex.faceRingGeneratorPowerCechExtDiagram_colimit_X_isZero, Reisner.Complex.faceRingGeneratorPowerCechExtDiagram_colimit_homology_isZero, Reisner.Complex.faceRingLocalCohomologyIsoGeneratorPowerCechExtZeroTermColimit}
\leanok
\end{theorem}
\begin{proof}
The localization transition maps are maps of $A$-algebras. Their induced
signed differential agrees with the existing coefficient-ring differential,
so the checked square-zero identity gives a complex in $A$-modules. The
fixed-stage Ext functor maps this complex to another complex. By
Theorem~\ref{thm:cech-acyclic-terms}, every positive-degree term is zero;
exactness and homology vanishing follow immediately.
In degree zero the Čech term is the face ring itself, so its Ext image is
canonically $\operatorname{Ext}^q_A(A/\mathfrak m^n,A)$. The vanishing of the next term
and the absent incoming differential at zero identify this with the degree-zero
homology of the Ext-applied complex.
\end{proof}

\begin{theorem}[Cech comparison: external foundation]\label{found:cech}
\uses{def:local-cohomology,thm:generator-power-ext-colimit,thm:localized-ext-rows,thm:cech-acyclic-terms}
Let $B$ be Noetherian and let $J=(f_v:v\in V)$ for a finite ordered set
$V$. For a $B$-module $M$, the extended cochain complex
\[
 C^q(f;M)=\bigoplus_{|F|=q}M_{f_F},\qquad f_F=\prod_{v\in F}f_v,
 \quad C^0=M,
\]
with alternating localization maps has cohomology naturally isomorphic to
the Ext-colimit model $H^q_J(M)$. For this project the required instance
is only $B=A=R[K]$, $M=A$, and $J=\mathfrak m=(\bar x_v:v\in V)$; use
this natural isomorphism to transport the Cech fine grading. No theorem
for arbitrary modules or ideals should be implemented first.
\end{theorem}
This follows from the Ext-colimit description in Stacks, Lemma 47.8.2,
and its extended Cech model in Lemma 47.9.1 and Section 47.10
\cite{stacks-cech}. Mathlib supplies the Ext-colimit definition but does
not supply this Cech comparison. Theorem~\ref{thm:ext-zero-cech-cohomology}
proves the degree-zero comparison for the face ring as an $R$-module.
Theorem~\ref{thm:generator-power-ext-colimit} supplies the generator-power
Ext stage system in all degrees, but does not identify it with Koszul
or Čech cohomology. In particular, the variables can be zero divisors,
so their Koszul complexes must not be assumed to resolve the quotients.
The grading-compatible comparison in all degrees remains an external
input to the informal proof. Inverting zero gives zero, which handles
ghost variables and nonface products.

\begin{theorem}[Depth detection: external foundation]\label{found:depth}
\uses{def:cm,def:local-cohomology}
If $B$ is Noetherian, $M$ is a finite $B$-module and $JM\ne M$, then
$\operatorname{depth}_J M=\min\{q\ge0:H^q_J(M)\ne0\}$, where
$\operatorname{depth}_J$ is maximal regular-sequence length in $J$.
For a local ring with $J$ maximal this is ordinary local depth.
\end{theorem}
This is Stacks, Lemma 47.11.1, tag 0AVZ \cite{stacks-depth}.
Mathlib supplies Ext-defined ideal depth and its regular-sequence
characterization. Theorem~\ref{thm:depth-cutoff} proves the finite-cutoff
form for a Noetherian local ring as a module over itself.

\begin{theorem}[Depth bound: external foundation]\label{found:depth-bound}
\uses{def:cm}
For a nonzero finite module $M$ over a Noetherian local ring $B$,
$\operatorname{depth} M\le\dim\operatorname{Supp} M$.
In particular $\operatorname{depth} B\le\dim B$.
\end{theorem}
This is Stacks, Lemma 10.72.3 \cite{stacks-depth-def}.

\begin{lemma}[Depth bound at the augmentation ideal]\label{lem:augmentation-depth-bound}
\uses{prop:dimension,def:cm}
Let $K$ be a simplicial complex on a finite ambient vertex set and $k$ a
field. For the localization $B=k[K]_{\mathfrak m}$ at the augmentation
ideal, $\operatorname{depth} B\le d(K)$, the maximum face cardinality.
\lean{Reisner.Complex.faceRing_atAugmentation_depth_le}
\leanok
\end{lemma}
\begin{proof}
\leanok
The Noetherian local ring $B$ has depth equal to the supremum of the
lengths of its regular sequences in the maximal ideal. Take one such
sequence of length $r$ and let $J$ be the ideal it generates.
The quotient $B/J$ is nonzero by regularity, so $\dim(B/J)\ge 0$.
The regular-sequence dimension formula gives $\dim(B/J)+r=\dim B$.
Thus $r\le \dim B=d(K)$ by Proposition~\ref{prop:dimension}.
Taking the supremum proves the bound.
\end{proof}

\begin{corollary}[Cohen--Macaulayness at the augmentation ideal]\label{cor:augmentation-local-cm}
\uses{lem:augmentation-depth-bound,prop:dimension,thm:depth-cutoff}
Let $B=k[K]_{\mathfrak m}$ be the augmentation localization of the face
ring of a simplicial complex on a finite ambient vertex set over a field
$k$, and let $\mathfrak n$ be the maximal ideal of $B$. Then
\[
  \operatorname{depth}B=\dim B
  \quad\Longleftrightarrow\quad
  H^q_{\mathfrak n}(B)=0\quad\text{for all }0\le q<d(K).
\]
\lean{Reisner.Complex.faceRing_atAugmentation_depth_eq_iff_localCohomology_vanishes}
\leanok
\end{corollary}
\begin{proof}
\leanok
By Theorem~\ref{thm:depth-cutoff}, the vanishing condition is equivalent
to $d(K)\le\operatorname{depth}B$. Lemma~\ref{lem:augmentation-depth-bound}
gives the reverse inequality and Proposition~\ref{prop:dimension}
identifies $d(K)$ with $\dim B$. When $K=\{\varnothing\}$, the cutoff
is zero, so the vanishing condition is vacuous and depth and dimension
are both zero.
\end{proof}

\begin{corollary}[Vanishing and the homogeneous maximal ideal test]\label{cor:homogeneous-maximal-cm-cutoff}
\uses{def:cm,lem:face-ring-graded-maximal,cor:augmentation-local-cm}
Let $A=k[K]$ over a field, with finite ambient vertex set. Write
$B=A_{\mathfrak m}$, $\mathfrak n=\mathfrak m B$, and $d=d(K)$.
The following conditions are equivalent.
\begin{enumerate}
\item $A_{\mathfrak p}$ is a Cohen--Macaulay local ring for every
prime $\mathfrak p$ maximal among proper homogeneous ideals of $A$.
\item $H^q_{\mathfrak n}(B)=0$ for every $0\le q<d$.
\end{enumerate}
\lean{Reisner.Complex.faceRing_homogeneousMaximal_cm_iff_localCohomology_vanishes}
\leanok
\end{corollary}
\begin{proof}
\leanok
Lemma~\ref{lem:face-ring-graded-maximal} identifies the only ideal to
test with $\mathfrak m$, which is prime since $A/\mathfrak m=k$.
Apply Corollary~\ref{cor:augmentation-local-cm} to $B$.
\end{proof}

\begin{corollary}[Global Cohen--Macaulayness implies augmentation-local vanishing]\label{cor:global-cm-implies-vanishing}
\uses{def:cm,prop:augmentation,cor:augmentation-local-cm}
Under the same assumptions on $A$ and $B$, if $A$ is Cohen--Macaulay
then $H^q_{\mathfrak n}(B)=0$ for every $0\le q<d(K)$.
\lean{Reisner.Complex.faceRing_localCohomology_vanishes_of_isCohenMacaulay}
\leanok
\end{corollary}
\begin{proof}
\leanok
Apply the global algebraic definition at the augmentation prime
and then Corollary~\ref{cor:augmentation-local-cm}.
\end{proof}

\begin{corollary}[Cohen--Macaulayness detected by local cohomology]\label{cor:cm-lc}
\uses{found:graded-cm,cor:homogeneous-maximal-cm-cutoff,cor:global-cm-implies-vanishing}
Let $B=A_{\mathfrak m}$ and $\mathfrak n=\mathfrak mB$. Writing
$d=d(K)=\dim B$, the algebraic ring $A$ is Cohen--Macaulay if and only if
$H^q_{\mathfrak n}(B)=0$ for all $0\le q<d$.
\end{corollary}
\begin{proof}
The forward implication is Corollary~\ref{cor:global-cm-implies-vanishing}.
For the converse, Corollary~\ref{cor:homogeneous-maximal-cm-cutoff} gives
Cohen--Macaulayness at every maximal homogeneous ideal. The external
graded local test gives Cohen--Macaulayness of $A$.
\end{proof}

\section{Computing the Cech pieces}

\begin{lemma}[Localization kernels and transition maps]\label{lem:localization-kernel}
\uses{prop:basis,thm:membership}
\lean{Reisner.Complex.faceRingFaceProduct, Reisner.Complex.faceRingFaceProduct_eq_monomial, Reisner.Complex.faceRingFaceProduct_empty, Reisner.Complex.faceRingFaceProduct_dvd, Reisner.Complex.faceRingFaceProduct_sdiff, Reisner.Complex.FaceRingLocalization, Reisner.Complex.faceRingLocalizationFraction, Reisner.Complex.faceRingLocalizationFraction_zero_power, Reisner.Complex.faceRingLocalizationFraction_surjective, Reisner.Complex.monomial_mul_mem_stanleyReisnerIdeal_iff, Reisner.Complex.faceRingFaceProduct_pow_mul_mk, Reisner.Complex.faceRingFaceProduct_mul_mk_eq_zero_iff, Reisner.Complex.faceRingLocalization_algebraMap_mk_eq_zero_iff, Reisner.Complex.faceRingLocalization_algebraMap_eq_zero_iff, Reisner.Complex.faceRingLocalizationFraction_eq_zero_iff, Reisner.Complex.faceRingLocalization_monomialFraction_ne_zero_iff, Reisner.Complex.faceRingLocalization_nontrivial_iff, Reisner.Complex.faceRingLocalization_subsingleton_of_nonface, Reisner.Complex.faceRingLocalization_scalarMap_injective, Reisner.Complex.faceRingLocalizationEmptyEquiv, Reisner.Complex.faceRingLocalizationEmptyEquiv_apply, Reisner.Complex.faceRingLocalizationMap, Reisner.Complex.faceRingLocalizationMap_algebraMap, Reisner.Complex.faceRingLocalizationMap_fraction, Reisner.Complex.faceRingLocalizationMap_self, Reisner.Complex.faceRingLocalizationMap_comp}
\leanok
Let $R$ be a commutative ring, $A=R[K]$, and $F$ a finite set of ambient
vertices. Write $x_F=\prod_{v\in F}\bar x_v$ and let $A_{x_F}$ be the
actual ring localization at the powers of $x_F$, including localization
at zero. Every element has a polynomial fraction representation.
If $p$ is a polynomial and $r\geq0$, then
\[
 \frac{\bar p}{x_F^r}=0\text{ in }A_{x_F}
 \quad\Longleftrightarrow\quad x_F\bar p=0\text{ in }A
 \quad\Longleftrightarrow\quad
 \forall\alpha\in\operatorname{supp}p,\quad
       F\cup\operatorname{supp}\alpha\notin K.
\]
Equivalently, the kernel of $A\longrightarrow A_{x_F}$ is the annihilator
$\operatorname{Ann}_A(x_F)=\{a\in A:x_Fa=0\}$.
For nontrivial $R$, a monic monomial fraction
$\bar x^\alpha/x_F^r$ is nonzero exactly when
$F\cup\operatorname{supp}\alpha\in K$, and $A_{x_F}$ is nontrivial
exactly when $F\in K$. The localization for $F=\varnothing$ is
$R$-algebra-equivalent to $A$.
Whenever $F$ is a face, the scalar map $R\longrightarrow A_{x_F}$ is
injective, also when $R$ is trivial.

For $F\subseteq G$ there is a canonical $R$-algebra map
$\rho_{FG}:A_{x_F}\longrightarrow A_{x_G}$ taking the image of every
$a\in A$ to its image in the target. It takes a fraction to the same
fraction interpreted in the target, where $x_F$ is a unit. These maps are
identity maps for $F=G$ and compose as
$\rho_{GH}\rho_{FG}=\rho_{FH}$ for $F\subseteq G\subseteq H$.
Explicitly, using $x_G=x_{G\setminus F}x_F$,
\[
 \rho_{FG}\left(\frac{\bar p}{x_F^r}\right)
   =\frac{x_{G\setminus F}^{\,r}\bar p}{x_G^r}.
\]
All assertions allow infinite ambient vertex types and ghost vertices.
\end{lemma}
\begin{proof}
\leanok
Every denominator is a power of $x_F$, and surjectivity of the polynomial
quotient map supplies a polynomial numerator. Since the denominator is a
unit in the localization, $\bar p/x_F^r=0$ exactly when $\bar p/1=0$,
including at $r=0$. The general localization criterion gives
$\bar p/1=0$ exactly when $x_F^n\bar p=0$ for some $n\geq0$.

To see why one factor suffices here, lift $x_F^n$ to the monic polynomial
$x^{n\mathbf1_F}$. Multiplication translates each exponent $\alpha$ of $p$
to $\alpha+n\mathbf1_F$. Translation is injective and preserves the
coefficient, even over a ring with zero divisors. For every $n\geq1$,
\[
  \operatorname{supp}(\alpha+n\mathbf1_F)
     =\operatorname{supp}\alpha\cup F .
\]
Membership in the squarefree nonface ideal depends only on these supports,
by Proposition~\ref{prop:basis}. Thus $x_F^n\bar p=0$ for any positive
$n$ exactly when $x_F\bar p=0$, with the stated termwise test.
A witness $n=0$ means $\bar p=0$ and hence also $x_F\bar p=0$.

For comparison, over a field $k$ take $B=k[x,y]/(x^2y)$. Inverting $x$
kills $\bar y$ because $x^2\bar y=0$, yet $x\bar y\ne0$ because $xy$ is
outside the monomial ideal $(x^2y)$. Here membership distinguishes the
two exponents of $x$ even though the vertex supports agree.

A monic monomial has one nonzero term over a nontrivial coefficient ring,
which gives the stated nonvanishing criterion. Applying it to the constant
monomial shows that $1\ne0$ in the localization exactly when $F$ is a face.
For a nonface, $x_F=0$ already in $A$, and inverting it gives the zero ring.
The empty product is one, and localization at a unit gives the claimed
empty-set equivalence.
For scalar injectivity, a nonzero constant has exponent zero in its
polynomial support. Its image cannot vanish when $F$ is a face, by the
kernel criterion.

If $F\subseteq G$, the product $x_F$ divides $x_G$, hence its image is a
unit in $A_{x_G}$. The universal property of localization gives
$\rho_{FG}$ with the stated values on $A$. Algebra homomorphisms out of
a localization are determined by these values, proving the fraction rule,
identity, and composition assertions. The explicit denominator formula
follows by applying the map to the defining fraction equation and using
$x_G^r=x_{G\setminus F}^{\,r}x_F^r$; the target denominator is a unit, so
this equation determines the fraction uniquely.
\end{proof}

\begin{lemma}[Localized monomial pieces]\label{lem:localized}
\uses{prop:basis,prop:grading,lem:localization-kernel}
\lean{Reisner.Complex.natExponent, Reisner.Complex.natExponent_apply, Reisner.Complex.positiveSupport, Reisner.Complex.LocalizationAdmissible, Reisner.Complex.fractionDegree, Reisner.Complex.fractionDegree_apply, Reisner.Complex.localizationMonomial, Reisner.Complex.localizationMonomial_fraction, Reisner.Complex.localizationMonomial_eq_zero_of_not_admissible, Reisner.Complex.localizationDegrees, Reisner.Complex.localizationMonomialBasis, Reisner.Complex.localizationMonomialBasis_apply, Reisner.Complex.localizationMonomial_ne_zero_iff, Reisner.Complex.localizationMonomial_natExponent, Reisner.Complex.faceRingLocalizationComponent, Reisner.Complex.faceRingLocalizationComponent_eq_bot, Reisner.Complex.faceRingLocalizationComponent_finrank, Reisner.Complex.faceRingLocalization_isInternal, Reisner.Complex.faceRingLocalizationGradedAlgebra, Reisner.Complex.faceRingLocalizationMap_monomial, Reisner.Complex.faceRingLocalizationMap_component, Reisner.Complex.faceRingLocalizationMap_basis, Reisner.Complex.faceRingLocalizationFraction_mem_component, Reisner.Complex.faceRingLocalizationComponentBasis, Reisner.Complex.faceRingLocalizationComponentBasis_apply, Reisner.Complex.faceRingLocalizationMap_basis_eq_zero, Reisner.Complex.faceRingFineComponent_map_localization}
\leanok
Let $A=R[K]$ over a commutative ring $R$ and let $F\subseteq V$ be finite.
Write $\mathbb Z^{(V)}$ for the finitely supported integer exponents. For
finite $V$ this is $\mathbb Z^V$. For $a\in\mathbb Z^{(V)}$ put
$N(a)=\{v:a_v<0\}$ and $P(a)=\{v:a_v>0\}$.
Call $a$ admissible for $F$ when
\[
 N(a)\subseteq F\qquad\text{and}\qquad F\cup P(a)\in K.
\]
The localization $A_{x_F}$ has an internal $\mathbb Z^{(V)}$ grading in which
$\bar x^\alpha/x_F^r$ is homogeneous of degree $\alpha-r\mathbf1_F$,
where $\mathbf1_F$ denotes the indicator of $F$. The degree $a$ piece is
zero for inadmissible $a$. The localization has an $R$-module basis indexed
by its admissible degrees, whose vector in degree $a$ is
\[
 m_F(a)=\frac{\bar x^{\,a+r\mathbf1_F}}{x_F^r},
 \qquad r\in\mathbb N,\quad a+r\mathbf1_F\geq0 .
\]
This Laurent monomial spans its component. Over a field $k$ the component
is one dimensional with basis $m_F(a)$ exactly when $a$ is admissible.
For a nonnegative exponent $\alpha$, the component is the image of the
degree $\alpha$ component of $A$. For $F\subseteq G$ the map $\rho_{FG}$
preserves degrees. For each $a$ admissible for $F$, it sends $m_F(a)$ to
$m_G(a)$ when $a$ is admissible for $G$, and to zero otherwise.
\end{lemma}
\begin{proof}
\leanok
If $N(a)\subseteq F$, choose $r$ at least as large as the absolute value
of each negative coordinate. Then $\beta=a+r\mathbf1_F$ is nonnegative
and agrees with $a$ outside $F$, so
$F\cup\operatorname{supp}\beta=F\cup P(a)$. Increasing $r$ by $s$
multiplies numerator and denominator by $x_F^s$. Hence any two choices
give the same fraction, by comparison at their maximum.
For nontrivial $R$, Lemma~\ref{lem:localization-kernel} shows that this fraction is nonzero
exactly when $F\cup P(a)$ is a face.

Write any element as $\bar p/x_F^r$ and expand $p$ in monomials.
The term of exponent $\alpha$ has degree $a=\alpha-r\mathbf1_F$ and
$N(a)\subseteq F$. The kernel criterion kills the term when
$F\cup P(a)$ is a nonface, and every surviving term is its coefficient
times $m_F(a)$. This proves spanning. For independence, take a finite
relation among the $m_F(a)$ with distinct admissible indices.
Choose one $r$ large enough for every index in the relation. Its numerator
is $\sum_a c_a x^{a+r\mathbf1_F}$. The exponents are distinct, so each
nonzero $c_a$ is the coefficient of a term whose support together with $F$
is a face. The polynomial-fraction kernel criterion forces every $c_a$
to vanish. Thus the spans of the indicated vectors, with zero in the
other degrees, form an internal direct sum.

Products of monomial fractions add numerator exponents and denominator
powers. The product therefore lies in the indicated piece of degree
$a+b$, including when the numerator is killed by the kernel. The unit
has degree zero, as represented by $1/1$. This proves the algebra grading.
At $r=0$ the image of $\bar x^\alpha$ generates the degree $\alpha$
piece, including when this image is zero. Proposition~\ref{prop:grading}
identifies the fine component of $A$ with the span of $\bar x^\alpha$,
which proves the image assertion.
For a nonface $F$ every piece vanishes. For $F=\varnothing$ the admissible
degrees are precisely the nonnegative exponents with face support.
Their components agree with the fine components of
Proposition~\ref{prop:grading} under the empty-product equivalence.
These admissible components are nonzero when $R$ is nontrivial.
When $R$ is the zero ring, every component is zero.

Finally the fraction formula of Lemma~\ref{lem:localization-kernel}
multiplies the numerator by $x_{G\setminus F}^{\,r}$ when replacing
$x_F^r$ by $x_G^r$.
The new exponent is $a+r\mathbf1_G$. Its support together with $G$ is
$G\cup P(a)$, and $N(a)\subseteq F\subseteq G$. The same kernel criterion
therefore gives precisely the stated images of the basis vectors.
\end{proof}

\begin{lemma}[The ordered extended Cech complex]\label{lem:cech-complex}
\uses{lem:localization-kernel,lem:localized}
Let $V$ be a finite linearly ordered set, let $R$ be a commutative ring,
and set $A=R[K]$. For $q\geq0$ the extended Cech cochains are
\[
 C^q=\prod_{\substack{F\subseteq V\\ |F|=q}} A_{x_F}.
\]
This finite product is canonically the direct sum of the displayed
$R$-modules. The index runs over all subsets of the indicated size.
A nonface contributes the zero ring, including any subset containing a
ghost vertex. The unique summand in $C^0$ is $A_{x_\varnothing}\cong A$.
For $q>|V|$ the product has empty index set and is the zero module.

For a finite set $G$ and $v\in G$, put
$r_G(v)=|\{u\in G:u<v\}|$. For $z\in C^q$ and $|G|=q+1$ set
\[
 (d^q z)_G=\sum_{v\in G}(-1)^{r_G(v)}
    \rho_{G\setminus\{v\},\,G}\bigl(z_{G\setminus\{v\}}\bigr).
\]
These $R$-linear maps satisfy $d^{q+1}d^q=0$. For each finitely supported
integer degree $a$, the families with $z_F\in(A_{x_F})_a$ in every
coordinate form a subcomplex $C^\bullet_a$.
\lean{Reisner.Complex.CechIndex, Reisner.Complex.CechCochain, Reisner.Complex.CechCochainDirectSum, Reisner.Complex.cechCochainDirectSumEquiv, Reisner.Complex.cechDifferentialComponent, Reisner.Complex.cechDifferential, Reisner.Complex.cechDifferential_apply, Reisner.Complex.cechDifferential_comp_eq_zero, Reisner.Complex.cechComplex, Reisner.Complex.cechComplex_X, Reisner.Complex.cechComplex_d, Reisner.Complex.cechDegreeSubmodule, Reisner.Complex.cechDifferential_maps_degree, Reisner.Complex.cechDifferentialDegree}
\leanok
\end{lemma}
\begin{proof}
The summands, the canonical equivalence for the empty product of variables,
and the transition maps are supplied by Lemma~\ref{lem:localization-kernel}.
Each term in the differential is a coordinate projection followed by an
$R$-linear transition map and multiplication by a sign, so their sum is
$R$-linear. In particular at a singleton the differential from degree zero
is the localization map with sign $+1$.

Fix a target $H$ of size $q+2$. In the double sum for $(d^{q+1}d^qz)_H$,
each ordered pair of distinct vertices $u,v\in H$ starts at the same set
$H\setminus\{u,v\}$ as its reversed pair. Both compositions are the
canonical map from that summand to $A_{x_H}$. If $u<v$, erasing $u$
reduces the rank of $v$ by one, and erasing $v$ leaves the rank of $u$
unchanged. Thus the two total sign exponents differ by one. Pairing the
two orders cancels the entire sum, also in characteristic two.

By Lemma~\ref{lem:localized}, each transition preserves degree $a$.
The target component is an $R$-submodule, so it contains the signed sum.
The restricted differentials inherit the square-zero identity.
\end{proof}

\begin{lemma}[Cech differential on monomial cochains]\label{lem:cech-basis-action}
\uses{lem:cech-complex,lem:localized}
Let $V$ be finite and linearly ordered, let $R$ be a commutative ring,
and set $A=R[K]$. Fix a finitely supported integer degree
$a\in\mathbb Z^{(V)}$. For every $F\subseteq V$, define
$e_F(a)\in C^{|F|}$ to be the cochain supported at the $F$ coordinate,
with value $m_F(a)$ when $a$ is admissible in $A_{x_F}$ and value zero
otherwise. Thus inadmissible monomial coordinates, including nonface and
ghost-containing subsets, are zero. If $a$ is admissible at $F$, then
\[
 d^{|F|}e_F(a)=\sum_{v\in V\setminus F}
 (-1)^{|\{u\in F:u<v\}|}e_{F\cup\{v\}}(a).
\]
On the right, an inadmissible target contributes the zero cochain.
\lean{Reisner.Complex.cechInsertIndex, Reisner.Complex.cechMonomialCochain, Reisner.Complex.cechMonomialCochain_apply_self, Reisner.Complex.cechMonomialCochain_apply_of_ne, Reisner.Complex.cechDifferential_monomialCochain}
\leanok
\end{lemma}
\begin{proof}
Evaluate both sides at any coordinate $H$ with $|H|=|F|+1$. In the
coordinate formula for the differential, a term can be nonzero only when
$H\setminus\{v\}=F$, which determines $v$ uniquely and is equivalent to
$H=F\cup\{v\}$ with $v\notin F$. The differential sign at this coordinate
is $(-1)^{|\{u\in H:u<v\}|}=(-1)^{|\{u\in F:u<v\}|}$. By
Lemma~\ref{lem:localized}, the transition map sends $m_F(a)$ to $m_H(a)$ if
$a$ is admissible at $H$, and to zero otherwise. This proves the formula in
every coordinate, including inadmissible targets.
\end{proof}

\begin{lemma}[The termwise fine-degree decomposition of the Cech cochains]\label{lem:cech-termwise-decomposition}
\uses{lem:cech-complex,lem:localized}
Let $V$ be finite, let $R$ be a commutative ring, and let $K$ be a complex
containing the empty face, with ghost vertices allowed. For each $q\geq0$
and $a\in\mathbb Z^{(V)}$, let $C^q_a$ be the submodule of the extended
Cech cochains whose coordinate at every $F\subseteq V$ with $|F|=q$ lies
in the degree-$a$ component of $A_{x_F}$. Then the family of these
submodules is an internal direct sum:
\[
  \bigoplus_{a\in\mathbb Z^{(V)}} C^q_a\ \xrightarrow{\ \sim\ }\ C^q,
\]
where the map is the sum of the coordinatewise inclusions. Equivalently,
\[
  \operatorname{DirectSum.IsInternal}\bigl(a\mapsto C^q_a\bigr).
\]
The indexing retains every subset coordinate, including the empty subset,
nonfaces, and subsets containing ghost vertices. The statement holds over
every commutative ring, including the zero ring.
\lean{Reisner.Complex.cechCochainMonomialBasis, Reisner.Complex.cechDegreeSubmodule, Reisner.Complex.cechDegreeSubmodule_isInternal}
\leanok
\end{lemma}
\begin{proof}
For each localization coordinate, Lemma~\ref{lem:localized} provides the
basis of admissible Laurent monomials. The finite product of these bases
gives a basis of $C^q$ indexed by pairs $(F,a)$, where $|F|=q$ and $a$ is
admissible at $F$. In the Lean model this is
`Pi.basis (fun F => localizationMonomialBasis R F.1)`, indexed by
`Sigma F : CechIndex V q, localizationDegrees F.1`.

For fixed $a$, the submodule $C^q_a$ is the image under the inverse basis
coordinate map of the finitely supported coefficient functions supported
on the fiber of indices $(F,b)$ with $b=a$. Distinct degree fibers are
disjoint, so their submodules are independent. The union of all fibers is
the entire basis index type; hence the submodules span $C^q$. The standard
internal-direct-sum criterion (independence and spanning) gives the result.
The construction uses the actual localization at every subset. A nonface
has zero localization. At a face, a degree which is inadmissible
contributes the zero component in that degree.
\end{proof}

\begin{lemma}[Fine-degree components of Cech cocycles and boundaries]\label{lem:cech-cycle-boundary-components}
\uses{lem:cech-complex,lem:cech-termwise-decomposition}
Let $V$ be finite and linearly ordered, $R$ a commutative ring, and $K$
a complex on $V$ containing the empty face. For $q\geq 0$ and
$a\in\mathbb Z^{(V)}$, write $\pi_a^q:C^q\to C^q_a$ for the projection
from the internal direct sum of Lemma~\ref{lem:cech-termwise-decomposition},
and $d_a^q:C^q_a\to C^{q+1}_a$ for the restricted differential.
The projection recovers a cochain in $C^q_a$ and vanishes on $C^q_b$ for
$b\ne a$. Every cochain $z$ is the sum of its projections over the finite
set of degrees in its support. The projections satisfy
\[
  \pi_a^{q+1}(d^qz)=d_a^q(\pi_a^q z)
  \qquad(z\in C^q).
\]
In particular $d^qz=0$ if and only if $d_a^q(\pi_a^qz)=0$ for every $a$.
For $w\in C^{q+1}_a$, its image in $C^{q+1}$ lies in
$\operatorname{im}d^q$ if and only if $w\in\operatorname{im}d_a^q$.
This last assertion describes incoming boundaries in positive cochain
degrees. In cochain degree zero the incoming boundary module is zero.
\lean{Reisner.Complex.cechDecomposition, Reisner.Complex.cechDecomposeLinearEquiv, Reisner.Complex.cechHomogeneousComponent, Reisner.Complex.cechHomogeneousComponent_eq_self, Reisner.Complex.cechHomogeneousComponent_eq_zero_of_ne, Reisner.Complex.cechHomogeneousComponent_sum_support, Reisner.Complex.cechHomogeneousComponent_eq_zero_iff, Reisner.Complex.cechHomogeneousComponent_differential, Reisner.Complex.cechDifferential_eq_zero_iff_components, Reisner.Complex.cechBoundary_homogeneous_primitive}
\leanok
\end{lemma}
\begin{proof}
\leanok
Write $z=\sum_b z_b$ as its unique finite sum of homogeneous components.
Degree preservation gives $d^qz_b\in C^{q+1}_b$, so $d^qz=\sum_b d^qz_b$
is already the homogeneous decomposition of $d^qz$. Taking its component
of degree $a$ proves the displayed identity. An element of the direct sum
vanishes exactly when every component vanishes, which gives the cocycle
criterion. If $d^qy=w$ with $w$ homogeneous of degree $a$, projecting
gives $d_a^q(\pi_a^q y)=w$. Conversely a primitive in $C^q_a$ is also a
primitive after inclusion in $C^q$.
\end{proof}

\begin{lemma}[Cohomology of the fine-degree Cech subcomplexes]\label{lem:cech-cohomology-decomposition}
\uses{lem:cech-complex,lem:cech-termwise-decomposition,lem:cech-cycle-boundary-components}
Let $K$ be a complex on a finite linearly ordered set $V$ and let $R$
be a commutative ring. Write $C^\bullet$ for its extended Cech complex
and $C^\bullet_a$ for the complex of its degree-$a$ submodules with the
restricted differentials. Write $i_a:C^\bullet_a\to C^\bullet$ and
$p_a:C^\bullet\to C^\bullet_a$ for the inclusion and projection.
For every $q\geq 0$ these cochain maps induce an $R$-linear isomorphism
\[
 T:\bigoplus_{a\in\mathbb Z^{(V)}} H^q(C^\bullet_a)
    \xrightarrow{\ \sim\ } H^q(C^\bullet),\qquad
 T(\iota_a x)=H^q(i_a)(x),
\]
where $\iota_a$ inserts a class in the $a$ summand. The $a$ component of
$T^{-1}(h)$ is $H^q(p_a)(h)$. If $z$ is a cocycle representing $h$, the
nonzero components of $T^{-1}(h)$ occur only in the finite degree support
of $z$.
\lean{Reisner.Complex.cechComplexDegree, Reisner.Complex.cechComplexDegree_d, Reisner.Complex.cechDegreeInclusion, Reisner.Complex.cechDegreeProjection, Reisner.Complex.cechDegreeInclusion_projection, Reisner.Complex.cechDegreeInclusion_projection_of_ne, Reisner.Complex.cechCohomologyDegree, Reisner.Complex.cechCohomologyAssembly, Reisner.Complex.cechCohomologyAssembly_lof, Reisner.Complex.cechCohomologyProjection_assembly, Reisner.Complex.cechCohomologyAssembly_injective, Reisner.Complex.cechCohomologyAssembly_surjective, Reisner.Complex.cechCohomologyAssembly_bijective, Reisner.Complex.cechCohomologyDecomposition, Reisner.Complex.cechCohomologyDecomposition_lof, Reisner.Complex.cechCohomologyDecomposition_symm_apply, Reisner.Complex.cechCohomologyDecomposition_support_subset, Reisner.Complex.cechCohomologyQuotientEquiv, Reisner.Complex.cechCohomologyDegreeQuotientEquiv, Reisner.Complex.cechCohomology_outgoing, Reisner.Complex.cechCohomology_incoming, Reisner.Complex.cechCohomologyDegree_outgoing, Reisner.Complex.cechCohomologyDegree_incoming, Reisner.Complex.cechCohomologyQuotientEquiv_homologyπ, Reisner.Complex.cechCohomologyDegreeQuotientEquiv_homologyπ, Reisner.Complex.cechCohomology_boundaries_zero, Reisner.Complex.cechCohomologyDegree_boundaries_zero}
\leanok
\end{lemma}
\begin{proof}
\leanok
For either complex $D^\bullet$, let $Z^q=\ker d^q$. The identity $d^2=0$
puts the image of the preceding differential inside $Z^q$. Its image
there is $B^q$, with $B^0=0$ since the complex starts in degree zero.
The module-category homology of $D^\bullet$ identifies canonically with
$Z^q/B^q$. In mathlib's short-complex API, the comparison at $q$ is
\texttt{moduleCatHomologyIso}. Under this
identification the canonical map from cocycles sends $z$ to its class.

By Lemma~\ref{lem:cech-cycle-boundary-components}, $p_a$ commutes with the
differential. The inclusions commute by definition of the restricted
differential. Thus both maps take cycles to cycles and boundaries to
boundaries, so both act on cohomology. In particular $H^q(i_a)$ takes a
homogeneous cocycle class to the class of the same cocycle in $C^\bullet$.
The composition $p_ai_b$ is the identity for $a=b$ and zero otherwise.
The induced maps on cohomology satisfy these same identities. Applying
$H^q(p_a)$ to the image under $T$ of a finite family therefore recovers
its $a$ component.

Conversely, let $z$ represent $h\in H^q(C^\bullet)$. By
Lemma~\ref{lem:cech-termwise-decomposition}, $z$ is the finite sum of the
$i_a p_a z$ over its degree support. Every $p_a z$ is a cocycle.
Taking classes proves that $T$ sends the finite family
$(H^q(p_a)(h))_a$ to $h$. Both identities hold in degree zero as well.
\end{proof}

\begin{lemma}[Cech differential and decomposition]\label{lem:cech-grading}
\uses{lem:cech-complex,lem:localized,lem:cech-basis-action,lem:cech-termwise-decomposition,lem:cech-cycle-boundary-components,lem:cech-cohomology-decomposition}
Fix an order on $V$. For each $F$ admissible in degree $a$, let $e_F$ be
the cochain of degree $|F|$ with value $m_F(a)$ in coordinate $F$ and
zero elsewhere. In the Cech complex of the images of all variables, the
differential on this degree $a$ basis is
\[
 d e_F=\sum_{v\notin F}(-1)^{|\{u\in F:u<v\}|}e_{F\cup\{v\}},
\]
where inadmissible terms mean zero. The differential squares to zero,
preserves $a$, and cohomology decomposes as the direct sum of its degree
$a$ cohomologies. Explicitly, inclusions of the subcomplexes induce
\[
 \bigoplus_{a\in\mathbb Z^{(V)}} H^q(C^\bullet_a)
       \ \xrightarrow{\ \sim\ }\ H^q(C^\bullet).
\]
\lean{Reisner.Complex.cechMonomialCochain, Reisner.Complex.cechDifferential_monomialCochain, Reisner.Complex.cechDifferential_comp_eq_zero, Reisner.Complex.cechDifferential_maps_degree, Reisner.Complex.cechCohomologyDecomposition, Reisner.Complex.cechCohomologyDecomposition_lof}
\leanok
\end{lemma}
\begin{proof}
\leanok
At a target $G$ of size $|F|+1$, only a term with $G\setminus\{v\}=F$ can
contribute to $d e_F$. Such a vertex is unique, and $r_G(v)$ counts the
elements of $F$ below $v$. Lemma~\ref{lem:localized} sends $m_F(a)$ to
$m_G(a)$ for an admissible target and to zero otherwise. This gives the
displayed basis formula from the differential of Lemma~\ref{lem:cech-complex}.

In each cochain $z$, every coordinate has a finite homogeneous
decomposition. There are finitely many coordinates, so the union of their
degree supports is finite. Collecting equal degrees gives a unique finite
sum $z=\sum_a z_a$. If $dz=0$, degree preservation and uniqueness give
$dz_a=0$ for each $a$. Thus the map on cohomology is surjective.
For injectivity, suppose a finite sum of homogeneous cocycles satisfies
$\sum_a z_a=dy$. Decomposing $y=\sum_a y_a$ gives $z_a=dy_a$ by uniqueness,
so every class in that sum already vanishes in its own subcomplex. In
degree zero the incoming boundary group is zero, and the same argument
uses uniqueness alone.
\end{proof}

\begin{lemma}[Positive-coordinate contraction]\label{lem:contraction}
\uses{lem:cech-grading}
Let $V$ be finite and linearly ordered, $R$ a commutative ring, and
$a\in\mathbb Z^{(V)}$. If $P(a)\ne\varnothing$, the degree $a$ Cech complex
is contractible. More precisely, each $p\in P(a)$ gives a homotopy from
the identity of $C^\bullet_a$ to zero. Its maps
$h^{q+1}:C^{q+1}_a\to C^q_a$ satisfy $h^1d^0=1$ and
$d^{q-1}h^q+h^{q+1}d^q=1$ for $q>0$. Consequently
$H^q(C^\bullet_a)=0$ for all $q\geq0$.
\lean{Reisner.Complex.cechMonomialCochainDegree, Reisner.Complex.cechCochainMonomialBasis_apply, Reisner.Complex.cechMonomialCochainDegree_eq_zero, Reisner.Complex.cechDegree_linearMap_ext, Reisner.Complex.cechDifferentialDegree_monomial, Reisner.Complex.cechPositiveContraction, Reisner.Complex.localizationAdmissible_insert_positive, Reisner.Complex.localizationAdmissible_erase_positive, Reisner.Complex.cechContraction_sign_cancel, Reisner.Complex.cechPositiveContraction_monomial, Reisner.Complex.cechContractionMap_eq_id, Reisner.Complex.cechPositiveContractingHomotopy, Reisner.Complex.cechPositiveContractingHomotopy_hom, Reisner.Complex.cechPositiveContraction_smul_monomial, Reisner.Complex.cechPositiveContraction_zero, Reisner.Complex.cechPositiveContraction_succ, Reisner.Complex.cechPositiveSupport_contractible, Reisner.Complex.cechPositiveCohomology_isZero, Reisner.Complex.cechPositiveCohomology_subsingleton}
\leanok
\end{lemma}
\begin{proof}
\leanok
Choose $p\in P(a)$. Admissibility is invariant under inserting or removing
$p$: it does not affect $N(a)\subseteq F$, and $p$ is already in $P(a)$.
At an admissible coordinate the Laurent monomial spans the component
freely over $R$ by Lemma~\ref{lem:localized}. When $p\notin F$,
the transition from $F$ to $F\cup\{p\}$ sends this generator to the
corresponding generator. At an inadmissible coordinate both degree
components are zero. Thus the transition is invertible on these
components over any commutative ring. Define an $R$-linear map on
monomial cochains by $h(e_F)=0$ for $p\notin F$, and
\[
 h(e_F)=(-1)^{|\{u\in F:u<p\}|}e_{F\setminus\{p\}}
 \quad(p\in F).
\]
For $p\notin F$, the term inserting $p$ in $h d(e_F)$ is $e_F$;
all other terms vanish. For $p\in F$, the term reinserting $p$ in
$d h(e_F)$ is $e_F$. For any other inserted vertex $v$, its contributions
to $dh$ and $hd$ have opposite signs: if $v<p$ its insertion changes
the removal sign once; if $p<v$ removing $p$ changes the insertion sign
once. The admissible terms match in pairs, including when they are zero.
In Lean the same cancellation follows from $d^2=0$ as follows.
The cochain map $\Delta=dh+hd-1$ vanishes on the monomials missing $p$.
For an admissible $F$ containing $p$, put $E=F\setminus\{p\}$.
Then $\Delta d(e_E)=d\Delta(e_E)=0$. Every insertion into $E$ other
than $p$ still misses $p$. Thus only the signed term $\pm\Delta(e_F)$
can survive, which forces $\Delta(e_F)=0$. Linearity on the admissible
monomials gives $dh+hd=1$ in every positive cochain degree. At $q=0$ the
only possible coordinate is the empty set and only $h^1d^0$ contributes.
At the top degree the outgoing term is zero. Above it the cochain
module is zero. Taking cohomology, the identity induces the zero map
because the two cochain maps are homotopic.
\end{proof}

\begin{lemma}[Nonpositive Cech degrees and link faces]\label{lem:cech-nonpositive-indices}
\uses{def:link,lem:localized,lem:cech-complex,lem:cech-cohomology-decomposition}
\lean{Reisner.Complex.negativeSupport, Reisner.Complex.mem_negativeSupport, Reisner.Complex.negativeSupport_finite, Reisner.Complex.negativeSupport_subset_iff, Reisner.Complex.positiveSupport_eq_empty_of_nonpos, Reisner.Complex.localizationAdmissible_iff_of_nonpos, Reisner.Complex.nonpositiveLocalizationIndexEquivLink, Reisner.Complex.nonpositiveLocalizationIndexEquivLink_apply, Reisner.Complex.nonpositiveLocalizationIndexEquivLink_symm_apply, Reisner.Complex.nonpositiveLocalizationIndex_card, Reisner.Complex.nonpositiveCechIndexEquivLink, Reisner.Complex.nonpositiveCechIndexEquivLink_apply, Reisner.Complex.nonpositiveCechIndexEquivLink_symm_apply, Reisner.Complex.not_localizationAdmissible_of_card_lt_negativeSupport, Reisner.Complex.not_localizationAdmissible_of_negativeSupport_nonface, Reisner.Complex.cechDegreeSubmodule_eq_bot_of_no_admissible, Reisner.Complex.cechDegreeSubmodule_eq_bot_of_card_lt_negativeSupport, Reisner.Complex.cechDegreeSubmodule_eq_bot_of_nonpos_negativeSupport_nonface, Reisner.Complex.cechComplexDegree_isZero_of_nonpos_negativeSupport_nonface, Reisner.Complex.cechCohomologyDegree_isZero_of_nonpos_negativeSupport_nonface}
\leanok
For $a\in\mathbb Z^{(V)}$, its negative support $N=\{v:a_v<0\}$ is
finite. Suppose $a\leq0$ pointwise. Then $P(a)=\varnothing$, and a
finite set $F$ is admissible in degree $a$ if and only if
$N\subseteq F$ and $F\in K$.

If $N\in K$, the maps $F\mapsto F\setminus N$ and $T\mapsto N\cup T$
give inverse bijections between admissible sets and faces of
$\operatorname{lk}_K N$. For corresponding $F$ and $T$ we have
$F=N\sqcup T$ and $|F|=|N|+|T|$. Thus for each $q\geq0$ the same maps
give a bijection
\[
 \{F: |F|=q,\ F\text{ admissible in degree }a\}
 \;\longleftrightarrow\;
 \{T\in\operatorname{lk}_K N: |N|+|T|=q\}.
\]
The empty face of the link corresponds to $F=N$ in degree $q=|N|$.
There are no admissible indices of size $q<|N|$. If $N\notin K$, there
are no admissible indices of any size.

For finite $V$ and any commutative ring $R$, the degree submodule
$C^q_a$ is zero whenever there is no admissible index of size $q$.
In particular it vanishes for $q<|N|$. If $N\notin K$, every
term of the ordered degree complex $C^\bullet_a$ and each of its
cohomology modules $H^q(C^\bullet_a)$ are zero.
\end{lemma}
\begin{proof}
\leanok
The negative coordinates form a subset of the finite support of $a$.
Their inclusion in $F$ says exactly that $a_v\geq0$ for each $v\notin F$.
When $a\leq0$ there are no positive coordinates, so the other
admissibility condition $F\cup P(a)\in K$ reduces to $F\in K$.
For a face $F$ containing $N$, the difference $T=F\setminus N$ is
disjoint from $N$ and satisfies $N\cup T=F$, hence is in the link.
Conversely, for a link face $T$, the face $N\cup T$ contains $N$ and
removing $N$ recovers $T$. Counting the disjoint union gives the
cardinality formula. If $N$ is a nonface, downward closure prevents
any face from containing it.

Each coordinate of a homogeneous cochain at an inadmissible $F$ belongs
to the zero localization component by Lemma~\ref{lem:localized}.
When all coordinates in a given degree vanish, the cochain itself is zero.
A complex is exact at any zero term, so its cohomology there is zero,
including in cochain degree zero.
\end{proof}

\begin{lemma}[Signed coordinate comparison with link cochains]\label{lem:cech-link-coordinate-comparison}
\uses{lem:cech-complex,lem:cech-basis-action,lem:cech-termwise-decomposition,lem:cech-nonpositive-indices,lem:augmented-cochain-complex,lem:augmented-basis-insertion}
\lean{Reisner.Complex.cechLinkSign, Reisner.Complex.cechLinkSign_empty, Reisner.Complex.cechLinkSign_mul_self, Reisner.Complex.cechLinkSign_insert, Reisner.Complex.cechRankBelow_union, Reisner.Complex.cechLinkSign_insertion_compat, Reisner.Complex.nonpositiveCechLinkIndexEquiv, Reisner.Complex.nonpositiveCechLinkIndexEquiv_val, Reisner.Complex.nonpositiveCechLinkIndexEquiv_symm_val, Reisner.Complex.cechNonpositiveLinkEquiv, Reisner.Complex.cechNonpositiveLinkEquiv_apply, Reisner.Complex.cechNonpositiveLinkEquiv_monomial, Reisner.Complex.cechNonpositiveLinkEquiv_symm_basis, Reisner.Complex.cechNonpositiveLinkEquiv_symm_apply, Reisner.Complex.cechNonpositiveLinkEquiv_differential, Reisner.Complex.cechNonpositiveLinkMap_monomial_apply, Reisner.Complex.cechLinkInsert_iff, Reisner.Complex.insert_link_admissible_iff}
\leanok
Let $V$ be finite and linearly ordered, $R$ a commutative ring, and
$a\in\mathbb Z^{(V)}$ pointwise nonpositive. Suppose its negative support
$N=N(a)$ is a face of $K$, and put $L=\operatorname{lk}_K N$.
For $r\in\mathbb N$, set $q=|N|+r$. An admissible Cech index of size
$q$ has the unique form $F_T=N\sqcup T$, where $T\in L$ and $|T|=r$.
Write $c_{F_T}(a)\in C^q_a$ for its Laurent monomial cochain, and
$e_T\in E^r(L;R)$ for the link's coordinate basis cochain. Define
\[
 s_N(T)=\sum_{t\in T}|\{n\in N:n<t\}|,\qquad
 \sigma_N(T)=(-1)^{s_N(T)}\in R.
\]
There is an $R$-linear equivalence
\[
 \Phi_r:C^{|N|+r}_a\xrightarrow{\ \sim\ }E^r(L;R),
 \qquad
 \Phi_r(c_{F_T}(a))=\sigma_N(T)e_T.
\]
More explicitly, write $z=\sum_{T\in\mathcal F_r(L)}
b_Tc_{F_T}(a)$ using the actual degree-fiber Laurent basis. Then
\[
 (\Phi_r z)(T)=\sigma_N(T)b_T,\qquad
 \Phi_r^{-1}(f)=
 \sum_{T\in\mathcal F_r(L)}\sigma_N(T)f(T)c_{F_T}(a).
\]
These equivalences commute with the existing adjacent differentials:
\[
 \Phi_{r+1}\circ d_a^{|N|+r}
   =\delta_L^r\circ\Phi_r,\qquad r\in\mathbb N.
\]
The relation between cardinalities is $|N|+|T|=q$, with no truncated
subtraction. At $r=0$, the empty link face corresponds to $F=N$,
$\sigma_N(\varnothing)=1$, and the commuting equality includes the
dual augmentation. Inadmissible Cech coordinates, including ghost and
nonface coordinates, remain in the actual Cech product and vanish in
this homogeneous degree. All assertions include $N=\varnothing$,
the only-empty-face link, characteristic two, and the zero coefficient
ring. This is a termwise comparison with differential compatibility.
The shifted-complex and induced cohomology statements are formulated
separately in Lemma~\ref{lem:cech-link-cohomology-comparison}; comparison
with local cohomology remains a further claim.
\end{lemma}
\begin{proof}
\leanok
The actual product Laurent basis of the Cech cochain module is indexed
by pairs consisting of a subset coordinate $F$ and an admissible Laurent
degree at $F$. Restricting its coordinate representation to the degree
fiber $a$ gives a basis of the actual submodule $C^q_a$, indexed by the
admissible subsets of cardinality $q$. Apply the checked bijection of
Lemma~\ref{lem:cech-nonpositive-indices}, and cancel $|N|$ from
$|N|+|T|=|N|+r$, to reindex this basis by $\mathcal F_r(L)$.
The coefficient representation gives an unsigned linear equivalence
with $E^r(L;R)$. Multiply its $T$ coordinate by $\sigma_N(T)$.
Since $\sigma_N(T)^2=1$, this multiplication is its own inverse,
giving the displayed equivalence and inverse without any assumption
that $R$ is nontrivial.

It suffices to check the differential equality on an admissible
monomial cochain $c_{N\cup T}(a)$. The checked Cech insertion formula
sums over vertices $v\notin N\cup T$. For such a vertex,
$\operatorname{insert}(v,N\cup T)$ is admissible in degree $a$ exactly
when $\operatorname{insert}(v,T)$ is a face of $L$: disjointness from
$N$ is retained, and the remaining condition is membership of their
union in $K$. An invalid target monomial vanishes in its actual
localization, so it contributes zero. For valid targets, disjointness
of $N$ and $T$ gives
\[
 r_{N\cup T}(v)=r_N(v)+r_T(v),\qquad
 \sigma_N(T\cup\{v\})=\sigma_N(T)(-1)^{r_N(v)}.
\]
The first identity counts the two disjoint parts below $v$; the second
adds the new summand in the exponent. Consequently
\[
 (-1)^{r_{N\cup T}(v)}\sigma_N(T\cup\{v\})
       =\sigma_N(T)(-1)^{r_T(v)},
\]
because the two factors $(-1)^{r_N(v)}$ multiply to one. This is exactly
the coefficient in the link's checked insertion formula. Checking
target coordinates therefore proves the commuting equality, including
$T=\varnothing$. For $N=\varnothing$ the sign is identically one.
No division or cancellation by a nonzero scalar is used, so the argument
also applies in characteristic two and to the zero ring.
The Lean equivalence extracts the full Laurent-basis coefficients and
constructs the inverse by linear extension on the link basis. The checked
monomial-generation principle then verifies the inverse law and the
coordinatewise differential equality.
\end{proof}

\begin{lemma}[Shifted link complex and Cech cohomology comparison]\label{lem:cech-link-cohomology-comparison}
\uses{lem:cech-link-coordinate-comparison,lem:cech-nonpositive-indices,lem:augmented-cochain-complex,lem:augmented-cochains}
\lean{Reisner.Complex.cechLinkDegreeEmbedding, Reisner.Complex.cechLinkShiftedComplex, Reisner.Complex.cechLinkShiftedComplexXIso, Reisner.Complex.cechLinkShiftedComplex_d, Reisner.Complex.cechLinkShiftedComplex_isZero_of_lt, Reisner.Complex.cechNonpositiveLinkComplexComponentIso, Reisner.Complex.cechNonpositiveLinkComplexComponentIso_add, Reisner.Complex.cechNonpositiveLinkComplexIso, Reisner.Complex.cechNonpositiveLinkHomologyIso, Reisner.Complex.cechNonpositiveLinkNaturalHomologyIso, Reisner.Complex.cechNonpositiveLinkCohomologyIso, Reisner.Complex.cechCohomologyDegree_isZero_of_nonpos_card_lt_negativeSupport}
\leanok
Let $V$ be finite and linearly ordered, $R$ a commutative ring, and
$a\in\mathbb Z^{(V)}$ pointwise nonpositive. Suppose $N=N(a)$ is a
face, and set $L=\operatorname{lk}_K N$ and $m=|N|$.
Write $E_L^\bullet$ for the checked natural face-cardinality cochain
complex of $L$, and $\mathcal E_L^\bullet$ for its checked
integer-indexed extension, whose degree $i$ represents reduced cochain
degree $i$. The map
\[
 j_m:\mathbb N\longrightarrow\mathbb N,\qquad j_m(r)=m+r
\]
is an injective embedding of cochain shapes: it preserves the relation
from a degree to its successor. Extend $E_L^\bullet$ by zero along
$j_m$ to obtain an actual cochain complex $D^\bullet$ of $R$-modules.
It has canonical term identifications
\[
 D^q\cong0\quad(q<m),\qquad
 \rho_r:D^{m+r}\xrightarrow{\ \sim\ }E_L^r
       \quad(r\in\mathbb N).
\]
Under these identifications the differential from $m+r$ to $m+(r+1)$
is the existing $\delta_L^r$, with its existing orientation signs:
\[
 \rho_{r+1}\circ d_D^{m+r}=\delta_L^r\circ\rho_r.
\]
Every differential whose source is below $m$ is zero, including the
incoming map to $D^m$ when $m>0$.

There is an isomorphism of actual natural-indexed cochain complexes
\[
 \Psi:C^\bullet_a\xrightarrow{\ \sim\ }D^\bullet.
\]
At degree $m+r$, its component is $\rho_r^{-1}\Phi_r$, where
$\Phi_r$ is the signed coordinate equivalence of
Lemma~\ref{lem:cech-link-coordinate-comparison}; its inverse component
is $\Phi_r^{-1}\rho_r$. Below $m$ its components are the unique maps
between the corresponding zero objects.
Consequently the induced maps identify the actual homology objects
in every natural degree $q$, with inverse induced by $\Psi^{-1}$:
\[
 H^q(C^\bullet_a)\cong H^q(D^\bullet).
\]
For $q<m$ both objects are zero. For every $r\in\mathbb N$, there are
canonical identifications
\[
 H^{m+r}(C^\bullet_a)
   \cong H^r(E_L^\bullet)
   \cong H^{(r:\mathbb Z)-1}(\mathcal E_L^\bullet).
\]
In particular $r=0$ reaches integer degree $-1$ and retains the empty
link face and its checked cohomology endpoints. If $q\geq m$, the
unique natural $r$ with $q=m+r$ satisfies
$(q:\mathbb Z)-(m:\mathbb Z)-1=(r:\mathbb Z)-1$; the case $q<m$ is
handled separately by its zero object, without truncated subtraction.
These assertions include $N=\varnothing$, empty ambient types, the
only-empty-face link, ghost vertices, characteristic two, and the zero
coefficient ring. They concern the actual Cech and simplicial cochain
complexes; the comparison with local cohomology remains separate.
\end{lemma}
\begin{proof}
\leanok
Injectivity of $j_m$ is cancellation of $m$, and
$j_m(r+1)=j_m(r)+1$ proves preservation of the cochain relation.
Its image consists exactly of the natural degrees at least $m$.
Extension along this embedding supplies the stated terms, transported
differentials, and square-zero law, including the zero incoming map at
the cutoff.

For $q<m$, Lemma~\ref{lem:cech-nonpositive-indices} makes the actual
Cech degree module zero as well. Use the unique isomorphism between
these zero objects for the component of $\Psi$. Otherwise write
$q=m+r$, using $r=q-m$ only after proving $m\leq q$.
Transport the source and target objects along this cardinality equality,
and use the checked $\Phi_r$ followed by the inverse term identification
of the extension. On adjacent degrees at least $m$, the component
commuting square is exactly the already checked signed
differential-compatibility theorem. At a source degree below $m$,
both sides of the square are maps from a zero object and hence agree.
This includes the square crossing from $m-1$ to $m$ when $m>0$.
The inverse term maps give the inverse cochain map.

Apply the homology isomorphism induced by $\Psi$ in every degree.
Extension along $j_m$ identifies its homology at $m+r$ with the
homology of $E_L^\bullet$ at $r$, including $r=0$: both have zero
incoming boundary at the cutoff. Compose with the inverse of the
checked integer-extension homology isomorphism of
Lemma~\ref{lem:augmented-cochains} to obtain the displayed integer
degree. At $q<m$, the term of $D^\bullet$ is zero; exactness at a zero
term gives zero actual homology, and $\Psi$ transports this to the
Cech complex. The additive cardinality relation gives the stated
integer degree equality.

The complex isomorphism and its homology comparisons use the checked
component maps, the embedding extension, and mathlib homology transport.
The positive-coordinate and nonface-negative-support zero cases remain
the separate checked results. The positive-degree local-cohomology
comparison and the local-cohomology link formula remain future work; the
augmented chain complex is checked separately in
Lemma~\ref{lem:augmented-chain-complex}, while canonical homology duality
remains open.
\end{proof}

\begin{theorem}[Uniform vanishing for the actual Cech complex]\label{cor:cech-uniform-link}
\uses{lem:cech-grading,lem:contraction,lem:cech-nonpositive-indices,
lem:cech-link-cohomology-comparison}
Let $K$ be a complex on a finite linearly ordered ambient set $V$, let $R$
be a commutative ring, and let $d\in\mathbb N$. Then the actual extended
Cech complex satisfies
\[
 H^q(C^\bullet(K;R))=0\quad(0\le q<d)
\]
if and only if for every face $F\in K$ and every $r\in\mathbb N$ with
$|F|+r<d$, the augmented reduced cohomology of the link vanishes in degree
$r-1$:
\[
 \widetilde H^{r-1}(\operatorname{lk}_K F;R)=0.
\]
This statement is about the checked Cech complex and link cohomology. It does
not identify Cech cohomology with Ext-colimit local cohomology or replace
cohomology by homology.
\lean{Reisner.Complex.cechCohomology_vanishes_iff_linkCohomology_vanishes}
\leanok
\end{theorem}
\begin{proof}
\leanok
The checked fine-degree decomposition identifies total Cech cohomology with
the direct sum of its fine-degree cohomology objects. A strand with a positive
coordinate is contractible. In a nonpositive degree, a nonface negative
support gives the zero complex, and a face negative support $N$ gives zero
below $|N|$ and the reduced cohomology of $\operatorname{lk}_K N$ in degree
$q-|N|-1$ at or above the cutoff. This proves one direction by taking the
fine degree equal to $-1$ on a chosen face and zero elsewhere; the converse
applies the checked strand cases and the assumed vanishing for each face.
The argument includes the empty face and $r=0$, hence reduced degree $-1$.
When $d=0$ both sides are vacuous. No field or nonempty-vertex hypothesis is
needed.
\end{proof}

\begin{corollary}[Čech cutoff in link homology]\label{cor:cech-uniform-link-homology}
\uses{cor:cech-uniform-link,lem:dual}
Let $K$ be a finite augmented complex, $k$ a field, and $d\in\mathbb N$.
The actual extended Čech cohomology vanishes in degrees $q<d$ if and only if
for every face $F\in K$ and $r\in\mathbb N$ with $|F|+r<d$, the augmented
chain homology of $\operatorname{lk}_K F$ vanishes in face-cardinality
degree $r$ (reduced degree $r-1$). Equivalently, for every face $F$ and
every $i\in\mathbb Z$ with $|F|+i+1<d$, the integer reduced homology
$\widetilde H_i(\operatorname{lk}_K F;k)$ vanishes.
\lean{Reisner.Complex.cechCohomology_vanishes_iff_linkHomology_vanishes, Reisner.Complex.cechCohomology_vanishes_iff_linkHomologyInt_vanishes}
\leanok
\end{corollary}
\begin{proof}
The checked cohomological cutoff theorem reduces the claim to each link,
where Lemma~\ref{lem:dual} converts cohomology vanishing to homology
vanishing. For $i\geq-1$ put $r=(i+1).\mathrm{toNat}$; then $r-1=i$
and the two cutoffs agree. Degrees below $-1$ already vanish by the zero
extension. This includes $r=0$, where the empty face supplies the reduced
degree $-1$ term, and includes $d=0$ without imposing an extra condition.
\end{proof}

\begin{corollary}[Čech vanishing, purity, and link vanishing]\label{cor:cech-purity}
\uses{cor:cech-uniform-link,lem:link-calculus,def:dimension}
Let $K$ be a finite augmented complex and $k$ a field. Vanishing of the
actual extended Čech cohomology below $d(K)$ is equivalent to the conjunction
that every facet has cardinality $d(K)$ and that, for every face $F$,
\[
 \widetilde H^i(\operatorname{lk}_K F;k)=0
 \quad\text{for all }i<\dim(\operatorname{lk}_K F).
\]
This is a Čech-side characterization. It does not assert the missing
positive-degree comparison with algebraic local cohomology.
\lean{Reisner.Complex.cechCohomology_vanishes_iff_pure_and_linkCohomology_vanishes}
\leanok
\end{corollary}
\begin{proof}
The actual Čech cutoff gives uniform link-cohomology vanishing. A facet of
size less than $d(K)$ has link $\{\varnothing\}$, whose degree $-1$
cohomology is nonzero over a field; hence all facets have size $d(K)$.
Purity gives equality in the link-dimension bound. Conversely, under purity,
vanishing below each link's own dimension implies the uniform bound and thus
actual Čech vanishing. The empty face, facet links and reduced degree $-1$
are included.
\end{proof}

\begin{corollary}[Čech vanishing, purity, and reduced link homology]\label{cor:cech-purity-homology}
\uses{cor:cech-purity,lem:dual}
Let $K$ be a finite augmented complex and $k$ a field. The actual extended
Čech cohomology vanishes in every natural degree $q<d(K)$ if and only if
every facet has cardinality $d(K)$ and
\[
 \widetilde H_i(\operatorname{lk}_K F;k)=0
 \quad\text{for every face }F\text{ and every }i\in\mathbb Z
 \text{ with }i<\dim(\operatorname{lk}_K F).
\]
The purity conjunct and own-link dimension cutoffs are retained exactly.
No connectivity hypothesis is added. The statement includes
$K=\{\varnothing\}$, ghost vertices, and facet links of dimension $-1$.
It remains a Čech-side statement pending comparison with algebraic local
cohomology.
\lean{Reisner.Complex.cechCohomology_vanishes_iff_pure_and_linkHomology_vanishes}
\leanok
\end{corollary}
\begin{proof}
Apply Corollary~\ref{cor:cech-purity} and replace the cohomology vanishing
condition for each link in each integer degree by the equivalent homology
vanishing condition from Lemma~\ref{lem:dual}. This leaves purity and all
dimension bounds unchanged. Below degree $-1$ the condition is automatic;
at degree $-1$ it uses augmented homology and retains the empty-face term.
\end{proof}

\begin{theorem}[Local-cohomology link formula]\label{thm:lc-link}
\uses{def:homology,lem:augmented-cochains,found:cech,lem:cech-grading,lem:contraction,lem:cech-nonpositive-indices,lem:cech-link-coordinate-comparison,lem:cech-link-cohomology-comparison}
For $q\ge0$ and $a\in\mathbb Z^V$,
\[
 H^q_{\mathfrak m}(A)_a\cong
 \begin{cases}
 \widetilde H^{q-|N(a)|-1}(\operatorname{lk}_K N(a);k),
   & a\le0\text{ and }N(a)\in K,\\
 0,&\text{otherwise}.
 \end{cases}
\]
This is a vector-space isomorphism, not an assertion about the full module
structure. It is the multigraded local-cohomology formula; see the
Hilbert-series version in Welker, Theorem 5.1 \cite{welker}.
\end{theorem}
\begin{proof}
Positive coordinates give zero by the contraction. Now suppose $a\le0$.
If $N=N(a)$ is a nonface, no face $F$ can contain $N$ and every term is
zero. If $N$ is a face, admissible sets have the unique form
$F=N\sqcup T$ for $T\in\operatorname{lk}_K N$.
The Cech degree is $q=|N|+|T|$, whereas $T$ has augmented cochain degree
$|T|-1$. For an explicit cochain identification multiply the basis
cochain of $T$ by
\[
 (-1)^{\sum_{t\in T}|\{n\in N:n<t\}|}.
\]
Inserting $v$ changes this exponent by $|\{n\in N:n<v\}|$.
That is precisely the difference between the Cech insertion sign in $F$
and the simplicial insertion sign in $T$. Thus the map commutes with the
differentials and identifies complexes with degrees shifted by $|N|+1$.
This includes $T=\varnothing$, the cochain of degree $-1$; omitting it
would give the wrong formula even for a simplex. Take cohomology and
apply the Cech comparison.
\end{proof}

\begin{corollary}[Uniform vanishing bound]\label{cor:uniform}
\uses{thm:lc-link,lem:dual}
For $d=d(K)$, the following are equivalent:
\[
 \begin{split}
 &H^q_{\mathfrak m}(A)=0\quad(0\le q<d);\\
 &\widetilde H_i(\operatorname{lk}_K F;k)=0
   \quad(F\in K,\ i<d-|F|-1).
 \end{split}
\]
\end{corollary}
\begin{proof}
For each face $F$ choose $a=-\mathbf1_F$, so $N(a)=F$. For
$-1\le i<d-|F|-1$, put $q=i+|F|+1$, which satisfies $0\le q<d$.
The formula and duality give the second assertion from the first. Degrees
$i<-1$ vanish by convention. Conversely, for every $0\le q<d$ and
every $a$, positive coordinates or a nonface $N(a)$ give a zero piece.
Otherwise $i=q-|N(a)|-1<d-|N(a)|-1$ and the assumed bound gives zero.
The direct-sum decomposition then makes the entire module zero.
For $d=0$ both assertions are vacuous after the below-minus-one convention.
\end{proof}

\section{Purity and the complete criterion}

\begin{lemma}[Linkwise vanishing forces purity]\label{lem:purity}
\uses{lem:link-calculus,lem:dual,lem:homology-connectivity}
If $\widetilde H_i(\operatorname{lk}_K F;k)=0$ for every face $F$ and
every integer $i<\dim\operatorname{lk}_K F$, then $K$ is pure: all facets
have the same cardinality.
\end{lemma}
\begin{proof}
Induct on $\dim K\ge-1$. Dimension $-1$ has only the empty facet, and
dimension zero has only singleton facets. For dimension at least one,
the hypothesis for the empty face gives $\widetilde H_0(K;k)=0$, so the
vertex-edge graph is connected by Lemma~\ref{lem:homology-connectivity}.
For an actual vertex $v$, its link has
smaller dimension and satisfies the same hypothesis: every link inside it
is the link of a face containing $v$, by link calculus. The induction
hypothesis makes this vertex link pure. Its facets correspond to facets
of $K$ containing $v$, so all those facets have a common size $r(v)$.
Every vertex belongs to a facet by finiteness. If $\{v,w\}$ is an edge,
extend it to a facet; that facet shows $r(v)=r(w)$. Connectedness now
makes $r(v)$ independent of $v$. Every facet is nonempty in these
dimensions and contains a vertex, so all facets have that size. This
derives purity and does not add it as an extra assumption.
\end{proof}

\begin{theorem}[Reisner's criterion]\label{thm:reisner}
\uses{cor:cm-lc,cor:uniform,lem:purity,lem:link-calculus}
\lean{Reisner.Complex.LinkHomologyVanishesBelowDimension, Reisner.Complex.reisner_criterion}
Let $k$ be any field and $K$ a complex containing the empty face on a
finite ambient vertex set. Its Stanley--Reisner ring is Cohen--Macaulay, in the
algebraic sense of Definition~\ref{def:cm}, if and only if
\[
 \widetilde H_i(\operatorname{lk}_K\sigma;k)=0
 \quad\text{for every }\sigma\in K\text{ and every }
 i\in\mathbb Z\text{ with }i<\dim\operatorname{lk}_K\sigma.
\]
\end{theorem}
\begin{proof}
Write $d=d(K)$. If the ring is Cohen--Macaulay, Corollaries
\ref{cor:cm-lc} and \ref{cor:uniform} give vanishing for
$i<d-|\sigma|-1$. Link calculus gives
$\dim\operatorname{lk}_K\sigma\le d-|\sigma|-1$, so this includes the
required range. Conversely the displayed condition makes $K$ pure by
Lemma~\ref{lem:purity}. Link calculus then gives equality
$\dim\operatorname{lk}_K\sigma=d-|\sigma|-1$ for every face.
The condition is therefore the uniform bound; the two corollaries imply
Cohen--Macaulayness. If $K=\{\varnothing\}$, its ring is $k$ and $d=0$;
its only link has dimension $-1$ and every required group is below $-1$.
Thus this edge case is included without changing the statement.
\end{proof}

\section{Formalization boundary and next steps}
The implemented entry points are Theorem~\ref{thm:membership},
Proposition~\ref{prop:basis}, Lemma~\ref{lem:homogeneous}, and
Propositions~\ref{prop:grading}, \ref{prop:augmentation}, and
\ref{prop:facet-primes}, together with Proposition~\ref{prop:dimension}
and Lemmas~\ref{lem:localization-kernel}, \ref{lem:localized},
\ref{lem:cech-complex}, and \ref{lem:cech-grading}, with the
augmented complex, squarefree exponent, ideal, quotient, and indicator
comparison that they need. Polynomial membership, the genuine quotient
basis and its monomial identification, nontriviality, scalar injectivity,
and weighted homogeneity of the polynomial ideal are checked. The latter
includes preservation of weighted homogeneous components and the fine
and total weight identities. The actual quotient now carries checked
weighted graded-algebra structures, including its fine and total gradings,
their coarsening relation, and the scalar algebra equivalence in degree
zero. Its augmentation, positive-degree and variable-ideal identifications,
homogeneous maximal ideal over a field, and residue algebra are checked.
Finite type for finite ambient vertices and Noetherianity over a Noetherian
coefficient ring are also checked. The coordinate quotients now have checked
algebra equivalences and Krull dimensions $|F|$, and the finite facet
intersection and prime-containment arguments are checked. Maximum face
cardinality, simplicial dimension, and the global and augmentation-local
face-ring dimension equalities are also checked. The actual variable-product
localizations, polynomial fraction kernels, monomial nonvanishing criteria,
scalar injectivity, empty-set equivalence, and coherent transition maps are
checked. Lemma~\ref{lem:localized} supplies the integer grading and Laurent
basis of each localization. Over a field each component has dimension zero
or one, and the transition maps carry each basis vector to the corresponding
target vector or to zero according to target admissibility.

Proposition~\ref{prop:dimension} now proves
$\dim A=\dim A_{\mathfrak m}=d(K)$ using both prime-chain bounds and
the augmentation-ideal height computation. Lemma~\ref{lem:cech-complex}
assembles the localization maps into the ordered extended Cech complex
and restricts its differential to each integer degree.
Lemma~\ref{lem:cech-complex} and Lemma~\ref{lem:cech-grading} are checked in
Lean, including the signed monomial differential formula, its square-zero
law and degree preservation, and the inclusion-induced direct-sum
decomposition of actual Čech cohomology. Lemma~\ref{lem:contraction} now
contracts each degree complex with a positive coordinate and proves
vanishing of its Cech cohomology. Definition~\ref{def:link} and
Lemma~\ref{lem:cech-nonpositive-indices} identify the admissible indices
in each nonpositive degree with the link faces of the required cardinality
when $N(a)$ is a face. For a nonface $N(a)$ the degree complex has zero
cohomology. Lemma~\ref{lem:augmented-cochain-complex} now constructs the
actual augmented face-cardinality cochain complex, its coordinate basis
and deletion formula, and its square-zero differential. The empty-face
scalar equivalence, constant augmentation on actual singleton faces,
zero outgoing map when all vertices are ghosts, and zero modules above
$|V|$ are checked. Lemmas~\ref{lem:augmented-basis-insertion} and
\ref{lem:augmented-cohomology-minus-one} give the coordinate insertion
formula and cardinality-zero cohomology endpoints.
Lemma~\ref{lem:augmented-cochains} now gives the actual integer-indexed
extension, its restriction and cohomology comparisons, zero terms and
cohomology below degree $-1$ and above degree $|V|-1$, and the transported
degree-minus-one endpoints. Lemma~\ref{lem:cech-link-coordinate-comparison}
now identifies each nonpositive Cech degree term with the corresponding
link cochains by an explicit sign twist and proves compatibility with
adjacent differentials over any commutative ring.
Lemma~\ref{lem:cech-link-cohomology-comparison} now packages these maps as
an actual isomorphism to the link complex extended by zero below $|N|$.
It identifies actual Cech cohomology in degree $|N|+r$ with the link's
integer-indexed cohomology in degree $(r:\mathbb Z)-1$, including $r=0$,
and proves vanishing below $|N|$. The augmented chain model and canonical
homology--cohomology duality are now checked in every integer degree,
including degree $-1$ and the low and high zero extensions. The link/facet
calculus, the uniform Čech cutoff, and its characterization as explicit
purity together with reduced link homology vanishing below each link's own
dimension are checked as well. Order independence, the graph-connectivity
comparison, the positive-degree comparison with local cohomology in
Theorem~\ref{thm:lc-link}, and the graded-local transfer remain future Lean
work. The statement of Theorem~\ref{thm:reisner} is present in Lean. Its
proof remains open.

\begin{thebibliography}{99}
\bibitem{panov} T. Panov, \emph{Cohomology of face rings, and torus actions},
arXiv:math/0506526v3 (2007), Theorems 2.5 and 2.9;
LMS Lecture Note Series 347 (2008), 165--201.
\url{https://arxiv.org/html/math/0506526v3}.
\bibitem{bh} W. Bruns and J. Herzog, \emph{Cohen--Macaulay Rings},
revised edition, Cambridge Studies in Advanced Mathematics 39,
Cambridge University Press, 1998, item 2.1.27 (graded local test).
\bibitem{bh-errata} W. Bruns and J. Herzog, corrections to the revised
edition, 20 May 2024, p. 1, note to 2.1.29 confirming 2.1.27.
\url{https://www.home.uni-osnabrueck.de/wbruns/brunsw/cmcorr.pdf}.
\bibitem{welker} V. Welker, \emph{Which properties of Stanley--Reisner
rings and simplicial complexes are topological?}, Section 5, Theorems
5.1--5.2. Used to check the local-cohomology shift, not as a substitute
for the Cech calculation given above.
\url{https://www.mathematik.uni-marburg.de/~welker/preprints/top_survey.pdf}.
\bibitem{stacks-cm} The Stacks Project, Sections 10.103--10.104,
algebraic Cohen--Macaulay modules and rings.
\url{https://stacks.math.columbia.edu/tag/00N7}.
\bibitem{stacks-cm-localize} The Stacks Project, Lemma 10.104.5,
localization of a Cohen--Macaulay local ring, tag 00NB.
\url{https://stacks.math.columbia.edu/tag/00NB}.
\bibitem{stacks-dimension} The Stacks Project, Sections 10.31, 10.60,
10.114, Noetherian rings and dimension of finite type algebras over fields.
\url{https://stacks.math.columbia.edu/tag/00OO}.
\bibitem{stacks-cech} The Stacks Project, Lemma 47.9.1 and Section 47.10,
extended Cech model and Noetherian comparison; Lemma 47.8.2, the Ext-colimit
description of local cohomology.
\url{https://stacks.math.columbia.edu/tag/0952}.
\bibitem{stacks-depth} The Stacks Project, Lemma 47.11.1, depth and
first nonzero local cohomology.
\url{https://stacks.math.columbia.edu/tag/0AVZ}.
\bibitem{stacks-depth-def} The Stacks Project, Section 10.72,
regular sequences, depth, and the dimension bound.
\url{https://stacks.math.columbia.edu/tag/00LE}.
\end{thebibliography}
